\documentclass[11pt]{amsart}

\usepackage[T1]{fontenc}
\usepackage{amsmath,amssymb,amsthm,mathtools}
\usepackage{booktabs,longtable,array}
\usepackage[margin=1.15in]{geometry}
\usepackage{url}
\usepackage[colorlinks=true,linkcolor=blue,citecolor=blue,urlcolor=blue]{hyperref}

\def\UrlBreaks{\do\,\do\.\do\_}
\urlstyle{tt}

\newtheorem{theorem}{Theorem}[section]
\newtheorem{proposition}[theorem]{Proposition}
\newtheorem{lemma}[theorem]{Lemma}
\newtheorem{corollary}[theorem]{Corollary}
\theoremstyle{definition}
\newtheorem{definition}[theorem]{Definition}
\newtheorem{assumption}[theorem]{Standing assumptions}
\theoremstyle{remark}
\newtheorem{remark}[theorem]{Remark}
\newtheorem{example}[theorem]{Example}

% Name of the Lean declaration(s) that formalize a statement (comma-separated, no spaces).
\makeatletter
\newcommand{\lname}[1]{\texttt{\detokenize{#1}}}
\newcommand{\@leanlist}[1]{\def\@leansep{}\@for\@leanname:=#1\do{\@leansep
  \texttt{\expandafter\detokenize\expandafter{\@leanname}}\def\@leansep{, }}}
\newcommand{\lean}[1]{\par\smallskip\noindent
  {\footnotesize\normalfont\raggedright\textsf{Lean:} \@leanlist{#1}\par}}
\makeatother
\newcommand{\notformal}{\textup{(not formalized)}}

\newcommand{\R}{\mathbb{R}}
\newcommand{\C}{\mathbb{C}}
\newcommand{\Z}{\mathbb{Z}}
\newcommand{\T}{\mathbb{T}}
\newcommand{\E}{\mathbb{E}}
\newcommand{\Prob}{\mathbb{P}}
\newcommand{\SO}{\mathrm{SO}}
\newcommand{\Or}{\mathrm{O}}
\newcommand{\tr}{\operatorname{tr}}
\newcommand{\sym}{\operatorname{sym}}
\newcommand{\skw}{\operatorname{skew}}
\newcommand{\re}{\operatorname{Re}}
\newcommand{\one}{\mathbf{1}}
\newcommand{\fro}[1]{\lVert #1\rVert_{F}}
\newcommand{\ip}[2]{\langle #1,#2\rangle}
\newcommand{\gap}{\operatorname{gap}}
\newcommand{\transp}{^{\mathsf T}}

\title[MCMC on $\mathrm{SO}(n)$ via representation theory]{Markov Chain Monte Carlo on $\mathrm{SO}(n)$\\ via Representation Theory}
\author{Daniel Kenigsberg}
\author{Avi Ray}

\begin{document}

\begin{abstract}
How quickly a Markov chain Monte Carlo sampler converges to its stationary distribution is
governed by its spectral gap.  We study the manifold sampler of Zappa, Holmes-Cerfon and Goodman
for the uniform distribution on the rotation group $\SO(n)$.  We show that this sampler is a
random walk on the group whose steps are invariant under conjugation, so that its spectrum is
computed by the representation theory of $\SO(n)$.  As a warm-up we diagonalize a random walk on
the dihedral group $D_n$ and find its spectral gap exactly.  For $\SO(2)$ with Gaussian tangent
proposals of variance $\sigma^2$ we determine the spectral gap exactly: it equals
$\E\bigl[(1-\sqrt{1-a^2})\,\one_{|a|\le 1}\bigr]$ with $a\sim N(0,\sigma^2)$, which is
$\sigma^2/2$ up to an error of at most $\tfrac32\sigma^4$.  For general $n$ we prove that every
matrix entry is an eigenfunction with eigenvalue $\lambda=\E[\tr R]/n$, where $R$ is the group
step, and deduce that the spectral gap is at most $1-\lambda\le (n-1)\sigma^2$, with
$1-\lambda$ equal to the mean squared displacement of one step divided by $2n$.  For $\SO(3)$ we
prove that this bound is sharp: the spectral gap of every conjugation-invariant walk whose steps
have positive semidefinite symmetric part, in particular of the sampler, is exactly
$1-\lambda$.  The proof develops the harmonic analysis of $\mathrm{SU}(2)$ from first principles.
On $\SO(4)$ the same statement is false for general conjugation-invariant walks, by an explicit
example, but it holds for every $n$ when the step distribution has a density that decreases
with each rotation angle.  The sampler has such a step distribution, so for Gaussian proposals we
obtain the matching lower bound $\gap\ge1-\lambda$, and the spectral gap on $\SO(n)$ is exactly
$1-\lambda$ for every $n$ and $\sigma$; the proof integrates by parts on the group and compares
gradients of matrix coefficients on the maximal torus.  We
also construct the sampler inside the formal system: the projection step, its reverse check, the
convergence of the Newton iteration with frozen Jacobian, and the conjugation invariance of the
step law are theorems, not modelling assumptions.  Every numbered proposition, lemma, theorem
and corollary has been formalized and machine-checked in Lean~4 with Mathlib, except the theorem
for general $n$ and its corollary for the sampler, which use the Peter--Weyl theorem; the lemmas
that carry their computations are formalized.  The text of this paper, its proofs and the
formalization were generated by an AI system, starting from the abstract of an earlier project
of the authors; Section~\ref{sec:verif-intro} says how, and what has and has not been checked.
\end{abstract}

\maketitle

\section{Introduction}

\subsection{The problem}
Let $X_0,X_1,X_2,\dots$ be a Markov chain with transition operator $P$ and stationary
distribution $m$.  If $P$ is reversible, the distance to stationarity of every square-integrable
observable decays geometrically, at a rate fixed by the spectrum of $P$ on $L^2(m)$.  The
\emph{spectral gap}
\[
  \gap(P)\;=\;\inf\Bigl\{\frac{\ip{f}{(I-P)f}}{\ip{f}{f}} \;:\; f\in L^2(m),\ \textstyle\int f\,dm=0,\ f\ne 0\Bigr\}
\]
measures this rate: the relaxation time of the chain is $1/\gap(P)$.  Bounding the spectral gap
of a sampler is therefore the basic quantitative question about it.

This paper studies that question for the sampler of Zappa, Holmes-Cerfon and Goodman
\cite{ZHG} on the rotation group
\[
  \SO(n)=\{X\in\R^{n\times n} : X\transp X=1,\ \det X=1\},
\]
with the uniform (Haar) distribution as target.  The sampler works on any manifold cut out by
equality constraints.  From the current point $X$ it draws a Gaussian step $V$ in the tangent
space at $X$, then moves from $X+V$ back to the manifold along the normal space at $X$ by
solving a nonlinear equation, and finally applies a Metropolis--Hastings test which includes a
check that the reverse move is possible.

\subsection{Results}
Our starting point is the observation that for $\SO(n)$ this sampler is a \emph{random walk on
the group}: one step is $X\mapsto XR$ for a random rotation $R$ that does not depend on $X$,
and the law $\nu$ of $R$ is invariant under conjugation (Section~\ref{sec:son}).  The transition
operator is then a convolution operator, and its spectrum is described by the Fourier transform
of $\nu$ at the irreducible representations of the group.  We carry this programme out in
increasing generality.

\begin{itemize}
\item \textbf{Finite groups: $D_n$} (Section~\ref{sec:dihedral}).  For the walk on the dihedral
group that multiplies by $e$, $r^{\pm1}$, $s$ with probabilities $a$, $b/2$, $b/2$, $c$, the
eigenvalues are exactly $a+b\cos(2\pi k/n)\pm c$, and for $n\ge2$ the spectral gap is
$\min\{2c,\ b(1-\cos(2\pi/n))\}$ (Theorems~\ref{thm:dihedral-spectrum}
and~\ref{thm:dihedral-gap}).

\item \textbf{$\SO(2)$} (Section~\ref{sec:sotwo}).  A tangent step of signed length $a$ projects
to a rotation by $\arcsin a$ when $|a|\le 1$ and has no projection otherwise.  For Gaussian
proposals $a\sim N(0,\sigma^2)$ the spectral gap is exactly
\[
  \gamma(\sigma)=1-\lambda_1(\sigma)=\E\bigl[(1-\sqrt{1-a^2})\,\one_{|a|\le1}\bigr],
  \qquad \bigl|\gamma(\sigma)-\tfrac12\sigma^2\bigr|\le\tfrac32\sigma^4,
\]
and $\lambda_1(\sigma)$ is also the $L^2$ contraction coefficient, so that
$\lVert P^tf\rVert_2\le\lambda_1(\sigma)^t\lVert f\rVert_2$ for mean-zero $f$
(Theorems~\ref{thm:so2-gap} and~\ref{thm:so2-size}).  The key step is an inequality showing
that the first Fourier mode dominates all others for every symmetric proposal density that is
non-increasing on $[0,1]$ (Theorem~\ref{thm:unimodal}).  The analysis is done on the circle
$\R/2\pi\Z$ and transported to the matrix group (Theorem~\ref{thm:so2-matrix}).

\item \textbf{$\SO(n)$} (Section~\ref{sec:son}).  For every conjugation-invariant step
distribution $\nu$ the matrix entries $X\mapsto X_{ij}$ are eigenfunctions of the walk with
eigenvalue $\lambda(\nu)=\E[\tr R]/n$ (Theorem~\ref{thm:entries}), so
$\gap\le 1-\lambda(\nu)$ (Theorem~\ref{thm:gap-upper}).  Moreover
$1-\lambda(\nu)=\E\fro{1-R}^2/(2n)$, and for the sampler
\[
  \frac{\E\bigl[\fro{A}^2;\ \text{accepted}\bigr]}{2n}\;\le\;1-\lambda(\nu)\;\le\;\frac{\E\fro{A}^2}{n}=(n-1)\sigma^2,
\]
where $A$ is the skew-symmetric tangent proposal with entries of variance $\sigma^2$
(Proposition~\ref{prop:one-minus-lambda}, Theorem~\ref{thm:sampler}).  In particular the
relaxation time is at least $1/((n-1)\sigma^2)$, and linear observables such as $\tr(BX)$
decorrelate at exactly the rate $\lambda(\nu)$.

\item \textbf{The sampler, constructed} (Section~\ref{sec:construct}).  For a skew-symmetric
$A$ the projection step has a solution with positive semidefinite symmetric part if and only
if $1+A^2\succeq0$; it is then unique, equal to $R(A)=\sqrt{1+A^2}+A$, and it is the limit of
Newton's iteration with the Jacobian frozen at the base point, which converges exactly in this
case (Theorem~\ref{thm:principal-exists}).  A move of the sampler, including the reverse check
and the Metropolis--Hastings test, is $X\mapsto XR(A)$ (Theorem~\ref{thm:zhg-move}).  For
Gaussian $A$ the law of $R(A)$ is conjugation invariant, so all the results above apply to it
without further hypotheses (Theorem~\ref{thm:gauss-sampler}).

\item \textbf{$\SO(3)$} (Section~\ref{sec:sothree}).  Let $\nu$ be a conjugation-invariant step
distribution on $\SO(3)$ whose steps have positive semidefinite symmetric part, that is,
rotation angle at most $\pi/2$.  Then the spectral gap is exactly $1-\lambda(\nu)$ and
$\lambda(\nu)$ is the $L^2$ contraction coefficient (Theorem~\ref{thm:so3-gap}); in particular
this holds for the sampler (Corollary~\ref{cor:so3-gauss}).  Mathlib has neither the Peter--Weyl
theorem nor the representation theory of $\SO(3)$, so the proof builds what it needs by hand on
the double cover $\mathrm{SU}(2)$: the spaces spanned by $k$-th powers of linear forms in the
matrix entries are eigenspaces of every conjugation-invariant walk
(Theorem~\ref{thm:su2-schur}); the eigenvalues satisfy a three-term recurrence which gives the
Weyl character formula (Theorem~\ref{thm:su2-weyl}); and the spaces are mutually orthogonal and
span a dense subspace (Propositions~\ref{prop:su2-complete} and~\ref{prop:su2-l2}).  The analytic
input is the inequality $|\chi_l(\theta)|/(2l+1)\le\chi_1(\theta)/3$ for $l\ge 1$ and
$\theta\in[0,\pi/2]$ (Theorem~\ref{thm:so3-char}).

\item \textbf{$\SO(n)$, $n\ge4$} (Section~\ref{sec:son-gap}).  The $2\times2$ minors of $X$ are
eigenfunctions with eigenvalue $\lambda_2(\nu)=\E[R_{11}R_{22}-R_{12}R_{21}]$, so
$\gap\le1-\lambda_2(\nu)$ as well (Theorem~\ref{thm:second}).  For the uniform distribution on
the conjugacy class of the rotation by $\pi/2$ in two orthogonal planes of $\R^4$, whose steps
have symmetric part $0$, one finds $\lambda=0$ and $\lambda_2=\tfrac13$: the spectral gap is at
most $\tfrac23$, not $1-\lambda=1$ (Theorem~\ref{thm:so4}).  So the theorem for $\SO(3)$ does not
extend to $\SO(4)$ in the same generality, nor to any even $n$; for odd $n\ge5$ we do not know
whether it does.  It does extend to every $n$ under a different hypothesis: if the
step distribution is $p\,\delta_1$ plus a multiple of a density with respect to Haar measure that
is nonincreasing in each rotation angle on $[0,\pi]$, then the spectral gap is exactly
$1-\lambda(\nu)$ and $\lambda(\nu)$ is the $L^2$ contraction coefficient, for every $n\ge2$
(Theorem~\ref{thm:monotone}).  The step distribution of the sampler has this form: its density
is a Gaussian in $\skw R$ times the Jacobian of $R\mapsto\skw R$, both of which decrease with
the rotation angles (Lemma~\ref{lem:jacobian}).  So for Gaussian proposals the spectral gap on
$\SO(n)$ is exactly $1-\E[\tr R]/n$ for every $n$ and every $\sigma$
(Corollary~\ref{cor:son-gauss}).  The proof integrates by parts on the group and compares, on
the maximal torus, the gradient of a matrix coefficient with that of the trace
(Lemma~\ref{lem:grad}).
\end{itemize}

Together with the upper bound of Section~\ref{sec:son}, Corollary~\ref{cor:son-gauss} determines
the spectral gap of the sampler with Gaussian proposals exactly, for every $n$ and $\sigma$: it is
$1-\E[\tr R]/n=\E\fro{1-R}^2/(2n)$, the mean squared displacement of one step divided by $2n$.

\subsection{Verification}\label{sec:verif-intro}
This paper was generated by an AI system.  Its starting point was the abstract of a project of
the two authors, which treated the dihedral group and $\SO(2)$ and left $\SO(n)$ as ongoing
work.  From that abstract the text, the statements and their proofs, the Lean development and
the numerical checks were produced by two of Anthropic's Claude language models, Claude
Fable~5.2 and Claude Opus~5.5, run as agents in Claude Code over several sessions directed by
the first author.  The paper has not been peer reviewed.  What stands behind each statement is therefore the following.

Every numbered Proposition, Lemma, Theorem and Corollary below has been formalized in Lean~4
\cite{Lean4} against Mathlib \cite{Mathlib} and checked by the Lean kernel, with two exceptions
that are marked \notformal: Theorem~\ref{thm:monotone} and Corollary~\ref{cor:son-gauss}.  Their
proofs combine the formalized Lemmas~\ref{lem:grad} and~\ref{lem:jacobian} with standard facts
that are not in Mathlib, chiefly the Peter--Weyl theorem (Remark~\ref{rem:monotone}).  These two
proofs were written by one of the two models and reviewed by the other, and they have been
tested numerically (Remark~\ref{rem:numerics}); this is weaker evidence than a formal proof or
a referee's report.  The names
of the Lean declarations are printed after each formalized statement, and
Section~\ref{sec:formal} gives the complete correspondence.  The prose between the numbered statements motivates and interprets
them and recalls standard background; it is not part of the formal development.  Results that
we assert but have \emph{not} formalized are stated in Remarks and Examples marked \notformal.
There are four kinds: the relation between the chain constructed here and an implementation of
the sampler of \cite{ZHG} (Remarks~\ref{rem:so2-model} and~\ref{rem:model}), classical results
that we quote for orientation but do not use (Remarks~\ref{rem:selfadjoint}
and~\ref{rem:so3-classical}), numerical observations (Remarks~\ref{rem:so2-numerics}
and~\ref{rem:numerics}), and heuristics.  They are listed again, with
Theorem~\ref{thm:monotone} and Corollary~\ref{cor:son-gauss}, at the end of
Section~\ref{sec:formal}.

\subsection{Related work}
The use of representation theory to analyse random walks on groups goes back to Diaconis and
Shahshahani; see \cite{Diaconis}.  Random walks on $\SO(n)$ generated by rotations in random
planes were analysed by Rosenthal \cite{Rosenthal}, and random walks on $\Or(n)$ generated by
random reflections by Porod \cite{Porod}, both using characters of the orthogonal group.  The
sampler studied here is from \cite{ZHG}.

\section{Random walks on groups}\label{sec:walks}

\subsection{Finite groups}
Let $G$ be a finite group and let $\mu\colon G\to\Bbbk$ be a function with values in a
commutative ring $\Bbbk$ (step weights; a probability vector when $\Bbbk=\R$, $\mu\ge0$ and
$\sum_g\mu(g)=1$).  The \emph{right random walk} with step weights $\mu$ moves from $x$ to $xg$
with probability $\mu(g)$.  Its Markov operator acts on functions $f\colon G\to\Bbbk$ by
\[
  (P_\mu f)(x)=\sum_{g\in G}\mu(g)\,f(xg).
\]

\begin{proposition}[Basic properties]\label{prop:finite-basic}
Let $\mu,\nu,f,h\colon G\to\Bbbk$.
\begin{enumerate}
\item[(a)] If $\sum_g\mu(g)=1$ then $P_\mu$ fixes constant functions.
\item[(b)] $\sum_x(P_\mu f)(x)=\bigl(\sum_g\mu(g)\bigr)\sum_xf(x)$.  In particular the uniform
distribution is stationary when $\sum_g\mu(g)=1$.
\item[(c)] If $\mu(g^{-1})=\mu(g)$ for all $g$, then
$\sum_x(P_\mu f)(x)\,h(x)=\sum_xf(x)\,(P_\mu h)(x)$: the walk is reversible.
\item[(d)] $P_\mu P_\nu=P_{\mu*\nu}$, where $(\mu*\nu)(g)=\sum_h\mu(h)\,\nu(h^{-1}g)$.
\end{enumerate}
\lean{walkOp_const,sum_walkOp,sum_walkOp_mul,walkOp_walkOp}
\end{proposition}

\begin{proof}
(a) is immediate.  For (b) and (c), exchange the two sums and substitute $x\mapsto xg^{-1}$ in
the inner one; in (c) also substitute $g\mapsto g^{-1}$ and use the symmetry of $\mu$.  For (d),
$(P_\mu P_\nu f)(x)=\sum_{h,k}\mu(h)\nu(k)f(xhk)$; substitute $k=h^{-1}g$.
\end{proof}

A \emph{matrix representation} of $G$ is a homomorphism $\rho$ from $G$ to the multiplicative
monoid of $d\times d$ matrices over $\Bbbk$.  The \emph{Fourier transform} of $\mu$ at $\rho$ is
the matrix
\[
  \hat\mu(\rho)=\sum_{g\in G}\mu(g)\,\rho(g).
\]

\begin{proposition}[Fourier analysis of the walk]\label{prop:finite-fourier}
Let $\rho$ be a matrix representation of $G$.
\begin{enumerate}
\item[(a)] For all $i,j$, $P_\mu(\rho_{ij})=\bigl(\rho\,\hat\mu(\rho)\bigr)_{ij}
=\sum_k\hat\mu(\rho)_{kj}\,\rho_{ik}$.  Thus the span of the matrix coefficients of $\rho$ is
invariant under $P_\mu$.
\item[(b)] If $\hat\mu(\rho)v=cv$, then for each $i$ the function $x\mapsto(\rho(x)v)_i$
satisfies $P_\mu f=cf$.
\item[(c)] If $\mu$ is a class function, $\mu(hgh^{-1})=\mu(g)$, then $\hat\mu(\rho)$ commutes
with $\rho(h)$ for every $h$.
\item[(d)] $\widehat{\mu*\nu}(\rho)=\hat\mu(\rho)\,\hat\nu(\rho)$.
\end{enumerate}
\lean{walkOp_matrixCoeff,walkOp_eigenfunction,fourierTransform_comm_of_class,fourierTransform_conv}
\end{proposition}

\begin{proof}
(a) Since $\rho(xg)=\rho(x)\rho(g)$, we get
$\sum_g\mu(g)\rho(xg)_{ij}=(\rho(x)\sum_g\mu(g)\rho(g))_{ij}$.  (b) follows by applying (a) to
the vector $v$.  For (c), $\rho(h)\hat\mu(\rho)=\sum_g\mu(g)\rho(hg)$; substituting
$g\mapsto h^{-1}gh$ and using that $\mu$ is a class function gives $\hat\mu(\rho)\rho(h)$.  (d)
is the same computation as Proposition~\ref{prop:finite-basic}(d).
\end{proof}

Part (c) is where representation theory enters: by Schur's lemma, if $\rho$ is irreducible and
$\mu$ is a class function then $\hat\mu(\rho)$ is a scalar, and by (a) every matrix coefficient
of $\rho$ is an eigenfunction.  In the examples below we never need Schur's lemma in general: we
compute $\hat\mu(\rho)$ directly.

To pass from a list of eigenfunctions to the whole spectrum we use the following criterion.  A
family $(\varphi_j)_{j\in J}$ of functions $G\to\C$, indexed by a finite set $J$, is
\emph{complete with constant $N$} if
\begin{equation}\label{eq:complete}
  \sum_{j\in J}\varphi_j(x)\,\overline{\varphi_j(y)}=N\cdot\one_{x=y}\qquad\text{for all }x,y\in G.
\end{equation}

\begin{proposition}[Complete families of eigenfunctions]\label{prop:eigenbasis}
Let $\mu\colon G\to\C$, let $(\varphi_j)_{j\in J}$ be complete with constant $N>0$, and suppose
$P_\mu\varphi_j=\lambda_j\varphi_j$ with $\lambda_j\in\R$ for every $j$.
\begin{enumerate}
\item[(a)] Suppose $\varphi_{j_0}$ is the constant function $1$ and $\lambda_j\le\Lambda$ for
all $j\ne j_0$.  Then for every $f\colon G\to\C$ with $\sum_xf(x)=0$,
\[
  \re\sum_x\overline{f(x)}\,(P_\mu f)(x)\;\le\;\Lambda\sum_x|f(x)|^2 .
\]
\item[(b)] Suppose $\mu$ is real valued and symmetric.  If $P_\mu f=cf$ for some $c\in\C$ and
some $f\ne0$, then $c=\lambda_j$ for some $j$.
\end{enumerate}
\lean{poincare_of_eigenbasis,eigenvalue_mem_of_eigenbasis}
\end{proposition}

\begin{proof}
Put $c_j=\sum_y\overline{\varphi_j(y)}f(y)$.  Completeness gives the expansion
$Nf=\sum_jc_j\varphi_j$, hence $N\,P_\mu f=\sum_j\lambda_jc_j\varphi_j$, and pairing with $f$
gives $N\sum_x\overline{f}\,P_\mu f=\sum_j\lambda_j|c_j|^2$ and $N\sum_x|f|^2=\sum_j|c_j|^2$.
In (a), $c_{j_0}=\sum_xf(x)=0$, and the claim follows.  In (b), $P_\mu$ is Hermitian, so
$c\,c_j=\ip{\varphi_j}{P_\mu f}=\ip{P_\mu\varphi_j}{f}=\lambda_jc_j$ for every $j$; if
$c\ne\lambda_j$ for all $j$ then all $c_j$ vanish and $f=0$.
\end{proof}

\subsection{Compact groups}\label{sec:compact}
\begin{assumption}\label{ass:compact}
In this subsection $G$ is a group with a $\sigma$-algebra for which multiplication and inversion
are measurable, $m$ is a probability measure on $G$ that is invariant under right translations,
and $\nu$ is a probability measure on $G$, the \emph{step distribution}.
\end{assumption}

The case of interest is a compact group with its Haar probability measure
(Lemma~\ref{lem:haar}).  The Markov operator of the right random walk
$X_{t+1}=X_tR_{t+1}$, with $R_1,R_2,\dots$ independent with law $\nu$, is
\[
  (P_\nu f)(x)=\int_Gf(xg)\,d\nu(g).
\]
In the formalization $f$ is an honest function, the integral is the Bochner integral, and
statements about $L^2(m)$ are statements about functions with $\int\lVert f\rVert^2\,dm<\infty$.

\begin{proposition}[Stationarity, contraction, reversibility]\label{prop:compact-basic}
Under the standing assumptions:
\begin{enumerate}
\item[(a)] If $f\in L^1(m)$ takes values in a Banach space, then $P_\nu f\in L^1(m)$ and
$\int P_\nu f\,dm=\int f\,dm$.
\item[(b)] If $f\in L^2(m)$ takes values in a normed space, then $P_\nu f\in L^2(m)$ and
$\int\lVert P_\nu f\rVert^2\,dm\le\int\lVert f\rVert^2\,dm$.
\item[(c)] If $\nu$ is invariant under inversion and $f,h\in L^2(m)$ are real or complex valued,
then $\int(P_\nu f)\,h\,dm=\int f\,(P_\nu h)\,dm$.
\end{enumerate}
\lean{integrable_markovOp,integral_markovOp,memLp_markovOp_and_le,integral_markovOp_mul}
\end{proposition}

\begin{proof}
The map $T(x,g)=(xg,g)$ preserves $m\otimes\nu$ because $m$ is right invariant.  Hence
$(x,g)\mapsto f(xg)=(f\circ\mathrm{pr}_1)(T(x,g))$ is integrable on the product whenever $f$ is
integrable on $G$, and Fubini's theorem gives (a).  For (b), Jensen's inequality gives
$\lVert P_\nu f(x)\rVert^2\le P_\nu(\lVert f\rVert^2)(x)$ for almost every $x$; integrate and use
(a).  For (c), apply the measure-preserving map $(x,g)\mapsto(xg^{-1},g)$ to
$(x,g)\mapsto f(xg)h(x)$ to get $\int\!\!\int f(x)h(xg^{-1})\,d\nu(g)\,dm(x)$, and then use the
invariance of $\nu$ under $g\mapsto g^{-1}$.
\end{proof}

\begin{lemma}[Compact groups are unimodular]\label{lem:haar}
Let $G$ be a compact topological group with its Borel $\sigma$-algebra.  Every Haar measure on
$G$ is right invariant.
\lean{isMulRightInvariant_of_compact}
\end{lemma}

\begin{proof}
For $g\in G$ the image of a Haar measure $m$ under $x\mapsto xg$ is a left-invariant finite
measure, hence a multiple $c\,m$ by uniqueness of Haar measure; comparing total masses gives
$c=1$.
\end{proof}

For a matrix-valued function $\pi\colon G\to\Bbbk^{d\times d}$, $\Bbbk\in\{\R,\C\}$, with
$\nu$-integrable entries, define $\hat\nu(\pi)=\int\pi\,d\nu$ entrywise.

\begin{proposition}[Fourier analysis on compact groups]\label{prop:compact-fourier}
Let $\pi$ be a matrix representation of $G$ with $\nu$-integrable entries.
\begin{enumerate}
\item[(a)] $P_\nu(\pi_{ij})=\bigl(\pi\,\hat\nu(\pi)\bigr)_{ij}$ for all $i,j$.
\item[(b)] If $\hat\nu(\pi)v=cv$, then $x\mapsto(\pi(x)v)_i$ satisfies $P_\nu f=cf$ for each $i$.
\item[(c)] Let $\pi$ be any matrix-valued function with $\nu$-integrable entries.  If
$\varphi\colon G\to G$ preserves $\nu$ and $\pi(\varphi(g))=B\,\pi(g)\,B'$ for all $g$ and fixed
matrices $B,B'$, then $\hat\nu(\pi)=B\,\hat\nu(\pi)\,B'$.
\end{enumerate}
\lean{markovOp_matrixCoeff,markovOp_eigenfunction,fourierMat_eq_conj}
\end{proposition}

Part (c) replaces Schur's lemma in Section~\ref{sec:son}: symmetries of the step distribution
force $\hat\nu(\pi)$ to commute with enough matrices to be a scalar.

\begin{definition}[Poincar\'e inequality, spectral gap, contraction coefficient]\label{def:gap}
The walk satisfies a \emph{Poincar\'e inequality with constant $\gamma$} if for every real
$f\in L^2(m)$ with $\int f\,dm=0$,
\[
  \int f\,(P_\nu f)\,dm\;\le\;(1-\gamma)\int f^2\,dm .
\]
The \emph{spectral gap} is the largest such $\gamma$, when it exists.  A number $\rho$ is an
\emph{$L^2$ contraction bound} if $\int|P_\nu f|^2\,dm\le\rho^2\int|f|^2\,dm$ for every complex
$f\in L^2(m)$ with $\int f\,dm=0$.
\lean{PoincareIneq,IsSpectralGap,L2ContractionBound}
\end{definition}

The Poincar\'e inequality is the same as $\gamma\operatorname{Var}_m(f)\le\ip{f}{(I-P_\nu)f}$,
the usual definition through the Dirichlet form.

\begin{remark}\label{rem:selfadjoint}
\notformal\ When $\nu$ is invariant under inversion, $P_\nu$ is a self-adjoint contraction of
$L^2(m)$ by Proposition~\ref{prop:compact-basic}, and by the spectral theorem the spectral gap
is $1-\sup\operatorname{spec}(P_\nu|_{\one^\perp})$ and the least contraction bound is the
spectral radius of $P_\nu|_{\one^\perp}$.  We do not use the spectral theorem: all results below
are stated and proved with Definition~\ref{def:gap}.
\end{remark}

\begin{proposition}\label{prop:gap-tools}
Under the standing assumptions:
\begin{enumerate}
\item[(a)] If $f\in L^2(m)$ is real, $\int f\,dm=0$, $\int f^2\,dm>0$ and
$P_\nu f=\lambda f$, then every Poincar\'e constant satisfies $\gamma\le1-\lambda$.
\item[(b)] If $\rho\ge0$ is an $L^2$ contraction bound, then the Poincar\'e inequality holds
with $\gamma=1-\rho$.
\item[(c)] If $\rho$ is an $L^2$ contraction bound, then for every complex $f\in L^2(m)$ with
$\int f\,dm=0$ and every $t\ge0$, $P_\nu^tf$ has mean zero and
$\int|P_\nu^tf|^2\,dm\le\rho^{2t}\int|f|^2\,dm$.
\end{enumerate}
\lean{PoincareIneq.le_of_eigenfunction,L2ContractionBound.poincareIneq,AddL2ContractionBound.iterate}
\end{proposition}

\begin{proof}
(a) Test the inequality on $f$.  (b) By Cauchy--Schwarz,
$\int f\,P_\nu f\le\lVert f\rVert_2\lVert P_\nu f\rVert_2\le\rho\lVert f\rVert_2^2$.  (c) By
induction on $t$, using Proposition~\ref{prop:compact-basic}(a) to see that $P_\nu^tf$ still has
mean zero.
\end{proof}

All statements of this subsection that involve $P_\nu$ are also available for groups written
additively (they are generated from the multiplicative versions), which is how they are applied
to the circle in Section~\ref{sec:sotwo}; part (c) of Proposition~\ref{prop:gap-tools} is
formalized in additive notation only.

\section{Finite groups: the dihedral group}\label{sec:dihedral}

Let $n\ge1$ and let $D_n$ be the dihedral group of order $2n$, with elements $r^i$ and $sr^i$
for $i\in\Z/n$, where $r^n=s^2=1$ and $srs=r^{-1}$.  Fix real numbers $a,b,c$ and consider the
step weights
\[
  \mu=a\,\delta_e+\tfrac b2\,(\delta_r+\delta_{r^{-1}})+c\,\delta_s ,
\]
so that $(P_\mu f)(x)=a f(x)+\tfrac b2\bigl(f(xr)+f(xr^{-1})\bigr)+c\,f(xs)$.  When $a,b,c\ge0$
and $a+b+c=1$ this is the lazy walk that rotates by one notch with probability $b$ and reflects
with probability $c$.  Put $\omega=e^{2\pi i/n}$.

\begin{lemma}[Two-dimensional representations]\label{lem:dihedral-rep}
For each $k\in\Z/n$ the assignment
\[
  \rho_k(r^i)=\begin{pmatrix}\omega^{ki}&0\\0&\omega^{-ki}\end{pmatrix},\qquad
  \rho_k(sr^i)=\begin{pmatrix}0&\omega^{-ki}\\\omega^{ki}&0\end{pmatrix}
\]
is multiplicative: $\rho_k(e)=1$ and $\rho_k(xy)=\rho_k(x)\rho_k(y)$ for all $x,y\in D_n$.
\lean{dihedralRep,dihedralRep_r,dihedralRep_sr}
\end{lemma}

\begin{proposition}[Fourier transform of the step weights]\label{prop:dihedral-fourier}
For every $k\in\Z/n$,
\[
  \hat\mu(\rho_k)=\begin{pmatrix}a+b\cos(2\pi k/n)&c\\c&a+b\cos(2\pi k/n)\end{pmatrix}.
\]
\lean{fourierTransform_dihedralRep,psi_re}
\end{proposition}

The matrix $\hat\mu(\rho_k)$ has eigenvectors $(1,\pm1)$ with eigenvalues
$a+b\cos(2\pi k/n)\pm c$.  By Proposition~\ref{prop:finite-fourier}(b), the first coordinate of
$x\mapsto\rho_k(x)(1,\varepsilon)\transp$ is an eigenfunction of the walk.

\begin{theorem}[Spectrum of the walk on $D_n$]\label{thm:dihedral-spectrum}
For $k\in\Z/n$ and $\varepsilon\in\{+1,-1\}$ define a function $\varphi_{k,\varepsilon}$ from $D_n$ to $\C$ by
\[
  \varphi_{k,\varepsilon}(r^i)=\omega^{ki},\qquad\varphi_{k,\varepsilon}(sr^i)=\varepsilon\,\omega^{-ki},
  \qquad\text{and put}\qquad \lambda_{k,\varepsilon}=a+b\cos(2\pi k/n)+\varepsilon c .
\]
\begin{enumerate}
\item[(a)] $P_\mu\varphi_{k,\varepsilon}=\lambda_{k,\varepsilon}\,\varphi_{k,\varepsilon}$.
\item[(b)] The $2n$ functions $\varphi_{k,\varepsilon}$ are complete with constant $2n$ in the
sense of \eqref{eq:complete}.
\item[(c)] A complex number is an eigenvalue of $P_\mu$ if and only if it equals
$\lambda_{k,\varepsilon}$ for some $k$ and $\varepsilon$.
\end{enumerate}
\lean{dihedralEigFun_r,dihedralEigFun_sr,walkOp_dihedralEigFun,dihedral_complete,dihedral_spectrum}
\end{theorem}

\begin{proof}
(a) was explained above.  (b) reduces to $\sum_{k\in\Z/n}\omega^{km}=n\cdot\one_{m=0}$ together
with $\sum_\varepsilon\varepsilon=0$.  (c) follows from (a), (b) and
Proposition~\ref{prop:eigenbasis}(b), since $\mu$ is real and symmetric.
\end{proof}

\begin{theorem}[Spectral gap of the walk on $D_n$]\label{thm:dihedral-gap}
Assume $a+b+c=1$ and $b,c\ge0$, and put
\[
  \gamma_n=\min\bigl\{2c,\ b\,(1-\cos(2\pi/n))\bigr\}.
\]
Then for every $f\colon D_n\to\C$ with $\sum_xf(x)=0$,
\[
  \re\sum_x\overline{f(x)}\,(P_\mu f)(x)\;\le\;(1-\gamma_n)\sum_x|f(x)|^2 .
\]
If $n\ge2$ there is a nonzero $f$ with $\sum_xf(x)=0$ for which equality holds (for all real
$a,b,c$ with $a+b+c=1$).  Thus $\gamma_n$ is the spectral gap when $n\ge2$.
\lean{dihedral_poincare,dihedral_poincare_sharp}
\end{theorem}

\begin{proof}
Apply Proposition~\ref{prop:eigenbasis}(a) with $j_0=(0,+)$, for which
$\varphi_{0,+}\equiv1$ and $\lambda_{0,+}=1$.  The other eigenvalues are $\lambda_{0,-}=1-2c$
and, for $k\ne0$, $\lambda_{k,\varepsilon}\le a+b\cos(2\pi/n)+c=1-b(1-\cos(2\pi/n))$, because
$\cos(2\pi k/n)\le\cos(2\pi/n)$ for $k\not\equiv0$.  Equality is attained by $\varphi_{0,-}$ or
by $\varphi_{1,+}$, whichever has the larger eigenvalue.
\end{proof}

The restriction to $n\ge2$ in the last statement is needed.  For $n=1$ the group has two
elements and $r=e$; by Theorem~\ref{thm:dihedral-spectrum} the eigenvalues of the walk are $1$
and $1-2c$, so its spectral gap is $2c$, whereas $\gamma_1=\min\{2c,0\}$.

\begin{example}
\notformal\ For fixed $b,c>0$ and large $n$ the gap is $b(1-\cos(2\pi/n))\sim2\pi^2b/n^2$: the
walk needs of order $n^2$ steps to mix, as expected of a diffusion on a cycle of length $n$.
The slow mode $\varphi_{1,+}$ comes from the two-dimensional representation $\rho_1$.
\end{example}

\section{The circle: \texorpdfstring{$\SO(2)$}{SO(2)}}\label{sec:sotwo}

\subsection{The sampler on \texorpdfstring{$\SO(2)$}{SO(2)}}
Write $R_\theta=\begin{psmallmatrix}\cos\theta&-\sin\theta\\\sin\theta&\cos\theta\end{psmallmatrix}$
and $J=\begin{psmallmatrix}0&-1\\1&0\end{psmallmatrix}$.  The tangent space of $\SO(2)$ at $X$
is spanned by $XJ$, and the normal space is $\{XS:S\transp=S\}$ (see Section~\ref{sec:son} for
general $n$).  A proposal is $X+aXJ$ with $a\in\R$, and the sampler looks for a symmetric $S$
such that $Y=X+aXJ+XS$ is orthogonal.  For a square matrix $M$ write
$\sym M=\tfrac12(M+M\transp)$ and $\skw M=\tfrac12(M-M\transp)$.

\begin{proposition}[$\SO(2)$ and the projection step]\label{prop:so2-model}
The following hold.
\begin{enumerate}
\item[(a)] $R_\theta\in\SO(2)$ and $R_{\theta+\varphi}=R_\theta R_\varphi$; every element of
$\SO(2)$ is $R_\theta$ for some $\theta$; and $R_\theta=1$ if and only if $\theta\in2\pi\Z$.
\item[(b)] If an orthogonal $2\times2$ matrix has skew part $aJ$, then $|a|\le1$.
\item[(c)] If $|a|\le1$, then $R_{\arcsin a}$ has skew part $aJ$ and positive semidefinite
symmetric part.
\end{enumerate}
\lean{rotMat_mem,rotMat_add,exists_rotMat_eq,rotMat_eq_one_iff,abs_le_one_of_skewPart_eq,principal_step_two}
\end{proposition}

By (a), $\theta\mapsto R_\theta$ induces a group isomorphism from $\T=\R/2\pi\Z$ onto $\SO(2)$.
By Proposition~\ref{prop:structure}(a) below, $Y=X+aXJ+XS$ with $S$ symmetric is equivalent to
$\skw(X\transp Y)=aJ$.  So a projection exists only if $|a|\le1$, and then the solution with
positive semidefinite symmetric part, which is the one continuously connected to $Y=X$ at $a=0$
and is unique by Proposition~\ref{prop:principal}(b), is $Y=XR_{\arcsin a}$.

\begin{definition}[Step angle and step distribution]\label{def:step}
The \emph{step angle} is $\theta(a)=\arcsin a$ if $|a|\le1$ and $\theta(a)=0$ otherwise (the
proposal is rejected).  For a probability measure $\varrho$ on $\R$ (the law of the tangent
coordinate $a$), the \emph{step distribution} $\nu_\varrho$ is the law of $\theta(a)\bmod2\pi$
on $\T=\R/2\pi\Z$.  We write $\nu_\sigma$ when $\varrho=N(0,\sigma^2)$.
\lean{stepAngle,circleStep,gaussStep}
\end{definition}

\begin{remark}[The model on the circle]\label{rem:so2-model}
Under the isomorphism $\theta\mapsto R_\theta$, the sampler on $\SO(2)$ with tangent proposals
of law $\varrho$ and principal projection becomes the random walk on $\T$ with step
distribution $\nu_\varrho$ and Markov operator $(Pf)(x)=\int f(x+y)\,d\nu_\varrho(y)$, and the
Haar measure of $\SO(2)$ becomes the normalized Haar measure of $\T$.  The theorems of
Sections~\ref{sec:circle-fourier}--\ref{sec:so2-gauss} are stated and proved for the walk on
$\T$.  They are transported to the matrix group in Theorem~\ref{thm:so2-matrix}, and
Theorem~\ref{thm:gauss-sampler}(d) identifies the result with the case $n=2$ of the sampler
constructed in Section~\ref{sec:construct}; both are formalized.  \notformal\ One point of
normalization: in the usual convention, where the tangent step is an isotropic Gaussian of
scale $s$ for the Frobenius norm, $a\sim N(0,s^2/2)$ because $\fro{XJ}=\sqrt2$; so
$\sigma^2=s^2/2$.
\end{remark}

\subsection{Random walks on the circle are Fourier multipliers}\label{sec:circle-fourier}
Let $m$ be normalized Haar measure on $\T$, $e_k(x)=e^{ikx}$ for $k\in\Z$, and
$\hat f(k)=\int\overline{e_k}\,f\,dm$.  For a probability measure $\nu$ on $\T$ put
$\hat\nu(k)=\int e_k\,d\nu$.  The characters $e_k$ are the irreducible representations of $\T$,
and every measure on an abelian group is conjugation invariant.

\begin{theorem}[Spectral decomposition on the circle]\label{thm:circle}
Let $\nu$ be a probability measure on $\T$ and $P=P_\nu$.
\begin{enumerate}
\item[(a)] $Pe_k=\hat\nu(k)\,e_k$ for every $k\in\Z$.
\item[(b)] $\widehat{Pf}(k)=\hat\nu(k)\,\hat f(k)$ for every $f\in L^1(m)$ and $k\in\Z$.
\item[(c)] If $|\hat\nu(k)|\le\rho$ for all $k\ne0$, then $\rho$ is an $L^2$ contraction bound.
Conversely every contraction bound $\rho\ge0$ satisfies $|\hat\nu(k)|\le\rho$ for all $k\ne0$.
Hence if $|\hat\nu(k)|\le|\hat\nu(k_0)|$ for all $k\ne0$ and some $k_0\ne0$, then
$|\hat\nu(k_0)|$ is the least nonnegative contraction bound.
\item[(d)] If $|\hat\nu(k)|\le\rho$ for all $k\ne0$ and $\hat\nu(k_0)=\rho$ for some $k_0\ne0$,
then the spectral gap of the walk is $1-\rho$.
\end{enumerate}
\lean{addMarkovOp_fourier,fourierCoeff_addMarkovOp,circle_l2ContractionBound,circleChar_le_of_bound,circle_isLeast_contraction,circle_isSpectralGap}
\end{theorem}

\begin{proof}
(a) $\int e_k(x+y)\,d\nu(y)=e_k(x)\hat\nu(k)$.  (b) By Fubini and translation invariance of $m$,
\[
  \widehat{Pf}(k)=\int\!\!\int\overline{e_k(x)}f(x+y)\,dm(x)\,d\nu(y)
  =\int e_k(y)\,d\nu(y)\int\overline{e_k(x)}f(x)\,dm(x).
\]
(c) Let $f\in L^2(m)$ have mean zero, so $\hat f(0)=0$.  By (b) and Parseval's identity for $Pf$
and for $f$,
\[
  \int|Pf|^2\,dm=\sum_k|\hat\nu(k)|^2|\hat f(k)|^2\le\rho^2\sum_{k}|\hat f(k)|^2=\rho^2\int|f|^2\,dm .
\]
The converse follows by testing on $e_k$, which has mean zero for $k\ne0$.  (d) The Poincar\'e
inequality with constant $1-\rho$ follows from (c) and Proposition~\ref{prop:gap-tools}(b).  It
is optimal by Proposition~\ref{prop:gap-tools}(a) applied to the real eigenfunction
$\re e_{k_0}$, which has mean zero, $\int(\re e_{k_0})^2\,dm=\tfrac12$, and eigenvalue $\rho$
because $\hat\nu(k_0)$ is real.
\end{proof}

Theorem~\ref{thm:circle} is formalized for the circle $\R/T\Z$ of any circumference $T>0$.

\subsection{The first mode dominates}
To use Theorem~\ref{thm:circle} we must compare all the coefficients $\hat\nu(k)$.  For the
sampler, $\hat\nu_\varrho(k)=\E[e^{ik\theta(a)}]$.

\begin{lemma}\label{lem:sin}
For every integer $k\ge0$ and every real $t$, $|\sin(kt)|\le k\,|\sin t|$.
\lean{abs_sin_nat_mul_le}
\end{lemma}

\begin{proof}
Induction on $k$, using $\sin((k+1)t)=\sin(kt)\cos t+\cos(kt)\sin t$.
\end{proof}

\begin{theorem}[Unimodal proposals]\label{thm:unimodal}
Let $\phi\colon\R\to\R$ be differentiable with continuous derivative, nonnegative, even and
integrable, and suppose $\phi'\le0$ on $[0,1]$.  Then for every integer $k\ge1$ and every real
$\varepsilon$ with $|\varepsilon|\le1$,
\[
  \int_{\R}\bigl(\cos\theta(a)+\varepsilon\cos(k\,\theta(a))\bigr)\,\phi(a)\,da\;\ge\;0 .
\]
\lean{unimodal_integral_nonneg}
\end{theorem}

\begin{proof}
The integrand is even, so the integral is twice the integral over $(0,\infty)$.  On $(1,\infty)$
we have $\theta(a)=0$ and the integrand is $(1+\varepsilon)\phi(a)\ge0$.  On $[0,1]$ substitute
$a=\sin t$ with $t\in[0,\pi/2]$:
\[
  \int_0^1\bigl(\cos\arcsin a+\varepsilon\cos(k\arcsin a)\bigr)\phi(a)\,da
  =\int_0^{\pi/2}G'(t)\,w(t)\,dt,
\]
where $G(t)=\sin t+\varepsilon\,\sin(kt)/k$ and $w(t)=\phi(\sin t)\cos t$.  Integrate by parts.
The boundary terms vanish because $G(0)=0$ and $w(\pi/2)=0$, so the integral equals
$-\int_0^{\pi/2}G(t)\,w'(t)\,dt$.  By Lemma~\ref{lem:sin}, $G\ge0$ on $[0,\pi]$.  Moreover
$w'(t)=\phi'(\sin t)\cos^2t-\phi(\sin t)\sin t\le0$ on $[0,\pi/2]$.
\end{proof}

Taking $\varepsilon=\pm1$ shows $|\E[\cos(k\theta(a))]|\le\E[\cos\theta(a)]$ whenever $a$ has
such a density $\phi$.  The inequality is not pointwise: $\cos(4\theta)>\cos\theta$ near
$\theta=\pi/2$.

\subsection{The spectral gap for Gaussian proposals}\label{sec:so2-gauss}
For $\sigma>0$ let $a\sim N(0,\sigma^2)$ and define
\[
  \lambda_1(\sigma)=\E[\cos\theta(a)],\qquad\text{where}\quad
  \cos\theta(a)=\begin{cases}\sqrt{1-a^2}&\text{if }|a|\le1,\\ 1&\text{if }|a|>1.\end{cases}
\]

\begin{theorem}[Spectral gap of the sampler on $\SO(2)$]\label{thm:so2-gap}
Let $\sigma>0$ and let $P$ be the Markov operator of the walk on $\T$ with step distribution
$\nu_\sigma$.
\begin{enumerate}
\item[(a)] $\hat\nu_\sigma(k)=\E[\cos(k\,\theta(a))]$ is real, and
$|\hat\nu_\sigma(k)|\le\lambda_1(\sigma)=\hat\nu_\sigma(1)$ for every $k\ne0$.
\item[(b)] The spectral gap of the sampler is exactly $1-\lambda_1(\sigma)$.
\item[(c)] $\lambda_1(\sigma)$ is the least nonnegative $L^2$ contraction bound; in particular,
for every complex $f\in L^2(m)$ with $\int f\,dm=0$ and every $t\ge0$,
\[
  \int|P^tf|^2\,dm\;\le\;\lambda_1(\sigma)^{2t}\int|f|^2\,dm .
\]
\end{enumerate}
\lean{gaussLam,cos_stepAngle,circleChar_gaussStep,circleChar_gaussStep_one,gauss_char_le,gauss_isSpectralGap,gauss_isLeast_contraction,gauss_iterate_bound}
\end{theorem}

\begin{proof}
The Gaussian law is symmetric and $\theta$ is odd, so $\hat\nu_\sigma(k)$ is real and equals
$\E[\cos(k\theta(a))]$.  The Gaussian density satisfies the hypotheses of
Theorem~\ref{thm:unimodal}, which gives (a) for $k\ge1$, and
$\hat\nu_\sigma(-k)=\hat\nu_\sigma(k)$.  Parts (b) and (c) follow from
Theorem~\ref{thm:circle}(c),(d) with $k_0=1$ and from Proposition~\ref{prop:gap-tools}(c).
\end{proof}

\begin{theorem}[Size of the gap]\label{thm:so2-size}
Let $\sigma>0$, $a\sim N(0,\sigma^2)$ and $\gamma(\sigma)=1-\lambda_1(\sigma)$.  Then
\begin{align*}
  \gamma(\sigma)&=\E\bigl[(1-\sqrt{1-a^2})\,\one_{|a|\le1}\bigr], &
  \tfrac12\,\E\bigl[a^2\,\one_{|a|\le1}\bigr]&\le\gamma(\sigma)\le\sigma^2, &
  \bigl|\gamma(\sigma)-\tfrac12\sigma^2\bigr|&\le\tfrac32\,\sigma^4 .
\end{align*}
\lean{one_sub_gaussLam,one_sub_gaussLam_ge,one_sub_gaussLam_le,integral_pow_four_gaussianReal,abs_one_sub_gaussLam_sub_le}
\end{theorem}

\begin{proof}
The first formula is $1-\cos\theta(a)=1-\sqrt{1-a^2}$ for $|a|\le1$ and $=0$ otherwise.  For
$x\in[0,1]$, $\tfrac x2\le1-\sqrt{1-x}\le\min\{x,\tfrac x2+\tfrac{x^2}2\}$, which gives
$\tfrac12a^2-\tfrac12a^4\le1-\cos\theta(a)\le\min\{a^2,\tfrac12a^2+\tfrac12a^4\}$ for every real
$a$ (for $|a|>1$ the left side is negative).  Take expectations and use $\E a^2=\sigma^2$ and
$\E a^4=3\sigma^4$.
\end{proof}

So the relaxation time of the sampler on $\SO(2)$ is $2/\sigma^2$ to leading order for small
steps, and the convergence is purely geometric with ratio $\lambda_1(\sigma)$.

\begin{remark}\label{rem:so2-numerics}
\notformal\ The gap is not monotone in $\sigma$: large proposals are rejected.  Numerical
quadrature gives a maximum $\gamma\approx0.133$ near $\sigma\approx0.76$.  The proof of
Theorem~\ref{thm:so2-gap} also goes through if the projection is only attempted for
$|a|\le c$ with some $c\le1$, for instance because the nonlinear solver gives up.
\end{remark}

\subsection{Back to the matrix group}\label{sec:so2-matrix}
The results above concern the walk on $\T=\R/2\pi\Z$.  The next theorem moves them to the group
of rotation matrices, where the sampler lives.

\begin{theorem}[The sampler on the matrix group $\SO(2)$]\label{thm:so2-matrix}
Let $\Phi\colon\T\to\SO(2)$ be the map induced by $\theta\mapsto R_\theta$, and let $m_2$ be the
normalized Haar measure of $\SO(2)$.
\begin{enumerate}
\item[(a)] $\Phi$ is well defined and bijective, $\Phi(x+y)=\Phi(x)\Phi(y)$, $\Phi$ is a
homeomorphism, and the image of the normalized Haar measure of $\T$ under $\Phi$ is $m_2$.
\item[(b)] Let $\varphi$ be a measurable bijection with measurable inverse from an additive group
$A$ onto a group $G$ such that $\varphi(x+y)=\varphi(x)\varphi(y)$, and let $m$ and $\nu$ be
measures on $A$.  A real number is a Poincar\'e constant (respectively an $L^2$ contraction
bound, respectively the spectral gap) for the walk on $A$ with step distribution $\nu$ and
reference measure $m$ if and only if it is one for the walk on $G$ with step distribution
$\varphi_*\nu$ and reference measure $\varphi_*m$.
\item[(c)] Let $\sigma>0$, let $a\sim N(0,\sigma^2)$, and let $\nu$ be the law of the rotation
matrix $R_{\theta(a)}$.  Then $\nu=\Phi_*\nu_\sigma$; the spectral gap of the walk
$X\mapsto XR$ on $\SO(2)$ with step distribution $\nu$ and stationary measure $m_2$ is exactly
$1-\lambda_1(\sigma)$; and $\lambda_1(\sigma)$ is its least nonnegative $L^2$ contraction bound.
\end{enumerate}
\lean{circleToSO_add,circleToSO_bijective,circleHomeo,map_circleToSO_haar,poincareIneq_map_iff,l2ContractionBound_map_iff,isSpectralGap_map_iff,so2Step,so2Step_eq_map,so2_isSpectralGap,so2_isLeast_contraction}
\end{theorem}

\begin{proof}
(a) By Proposition~\ref{prop:so2-model}(a), $\theta\mapsto R_\theta$ is a surjective homomorphism
$\R\to\SO(2)$ with kernel $2\pi\Z$, so $\Phi$ is a bijective homomorphism.  It is continuous, and
a continuous bijection from a compact space to a Hausdorff space is a homeomorphism.  The image
of Haar measure is a left-invariant probability measure on $\SO(2)$, hence equal to $m_2$ by
uniqueness of Haar measure.  (b) Composition with $\varphi$ is a bijection between functions on
$G$ and on $A$ which preserves integrals, $L^2$ norms and membership in $L^2$, and
$(P_{\varphi_*\nu}F)\circ\varphi=P_\nu(F\circ\varphi)$ because $\varphi$ is a homomorphism.
(c) The identity $\nu=\Phi_*\nu_\sigma$ holds by definition of $\nu_\sigma$.  Apply (b) with
$\varphi=\Phi$ to Theorem~\ref{thm:so2-gap}(b),(c), using (a).
\end{proof}

\section{The rotation group \texorpdfstring{$\SO(n)$}{SO(n)}}\label{sec:son}

\subsection{The sampler is a random walk on the group}\label{sec:reduction}
The group $\Or(n)=\{X:X\transp X=1\}$ is the zero set of the $n(n+1)/2$ constraints
$(X\transp X-1)_{ij}=0$, $i\le j$, in $\R^{n\times n}$ with the Frobenius inner product, and
$\SO(n)$ is its identity component.  At $X\in\Or(n)$ the tangent space is
$\{XA:A\transp=-A\}$ and the normal space, spanned by the gradients of the constraints, is
$\{XS:S\transp=S\}$.  One move of the sampler \cite{ZHG} from $X$ consists of a tangent proposal
$V=XA$ with $A$ skew-symmetric and random, followed by a search for a normal correction $XS$
such that $Y=X+XA+XS$ lies on the manifold.

\begin{proposition}[Structure of a move]\label{prop:structure}
Let $X$ be an orthogonal $n\times n$ matrix and $A$ a skew-symmetric one.
\begin{enumerate}
\item[(a)] For any matrix $Y$: there is a symmetric $S$ with $Y=X+XA+XS$ if and only if
$\skw(X\transp Y)=A$.
\item[(b)] Let $R$ be orthogonal, $P=\sym R$ and $A=\skw R$.  Then $PA=AP$ and $P^2=1+A^2$.
Conversely, if $P$ is symmetric, $A$ is skew-symmetric, $PA=AP$ and $P^2=1+A^2$, then $P+A$ is
orthogonal, with symmetric part $P$ and skew part $A$.
\item[(c)] For any matrix $R$, $\skw((XR)\transp X)=-\skw R$ and
$\sym((XR)\transp X)=\sym R$.  If $Y$ is orthogonal, $\fro{Y(-A)}=\fro{XA}$.
\end{enumerate}
\lean{exists_normal_iff,symPart_mul_skewPart_comm,symPart_sq,add_mem_orthogonalGroup,symPart_add_of_symm_of_skew,skewPart_add_of_symm_of_skew,skewPart_reverse,symPart_reverse,reverse_proposal_frobSq}
\end{proposition}

\begin{proof}
(a) $X\transp(X+XA+XS)=1+A+S$ has skew part $A$; conversely take $S=X\transp Y-1-A$.  (b)
Expanding $R\transp R=1=RR\transp$ with $R=P+A$ and $R\transp=P-A$ gives
$P^2-A^2\pm(PA-AP)=1$.  (c) $(XR)\transp X=R\transp$, and left multiplication by an orthogonal
matrix preserves the Frobenius norm.
\end{proof}

By (a), the candidate moves from $X$ with tangent proposal $XA$ are exactly the points $Y=XR$
with $R$ orthogonal and $\skw R=A$.  By (c), the reverse move from $Y=XR$ to $X$ is the group
step $R\transp$; it uses the tangent proposal $-YA$, which has the same length as $XA$.  By (b),
given $A$ the possible symmetric parts are the symmetric square roots of $1+A^2$ that
commute with $A$.  The \emph{principal} one is the positive semidefinite root; it exists when
$1+A^2\succeq0$, that is when the operator norm of $A$ is at most $1$, and it is the solution
continuously connected to $R=1$ at $A=0$.

\begin{proposition}[The principal root]\label{prop:principal}
Let $R$ be an orthogonal $n\times n$ matrix.
\begin{enumerate}
\item[(a)] If $\sym R$ is positive semidefinite, then $\det R=1$.
\item[(b)] If $R'$ is orthogonal, $\skw R'=\skw R$, and $\sym R$ and $\sym R'$ are positive
semidefinite, then $R'=R$.
\item[(c)] For every matrix $Q$, $\skw(QRQ\transp)=Q(\skw R)Q\transp$ and
$\sym(QRQ\transp)=Q(\sym R)Q\transp$, and the latter is positive semidefinite if $\sym R$ is.
\item[(d)] $\fro{\skw R}^2\le n$.
\end{enumerate}
\lean{det_eq_one_of_posSemidef,principal_root_unique,skewPart_conj,symPart_conj,posSemidef_symPart_conj,frobSq_skewPart_le_card}
\end{proposition}

\begin{proof}
(a) $\det R=\pm1$.  If $\det R=-1$, then
$\det(1+R)=\det(R(1+R)\transp)=-\det(1+R)$, so there is $v\ne0$ with $Rv=-v$, and then
$v\transp(\sym R)v=v\transp Rv=-|v|^2<0$.  (b) Both symmetric parts are positive semidefinite
square roots of $1+(\skw R)^2$, and such a root is unique.  (c) is a direct computation.  (d)
$\fro{\skw R}^2=n-\tr(P^2)$ with $P=\sym R$, and $\tr(P^2)=\fro{P}^2\ge0$.
\end{proof}

\begin{proposition}[Size of a step]\label{prop:step-size}
Let $R$ be orthogonal, $P=\sym R$ and $A=\skw R$.  Then
\[
  \fro{1-R}^2=2\,(n-\tr R)=\fro{A}^2+\fro{1-P}^2\;\ge\;\fro{A}^2,
\]
and if $P$ is positive semidefinite then also $\fro{1-R}^2\le2\,\fro{A}^2$.  Consequently, for
the principal root, $\tfrac12\fro{A}^2\le n-\tr R\le\fro{A}^2$.
\lean{frobSq_one_sub,frobSq_one_sub_eq,frobSq_skewPart_le,frobSq_one_sub_le,card_sub_trace_bounds}
\end{proposition}

\begin{proof}
The identities follow from $R\transp R=1$, $\tr R=\tr P$ and $\fro{A}^2=-\tr(A^2)=n-\tr(P^2)$.
The last inequality is $\tr(P^2)\le\tr P$, which follows from the identity
\[
  P-P^2=\tfrac12\,(1-P)\,P\,(1-P)+\tfrac12\,A\,P\,A\transp ,
\]
valid because $PA=AP$ and $A^2=P^2-1$: both terms on the right are positive semidefinite when
$P$ is, so their traces are nonnegative.
\end{proof}

We can now say precisely which Markov chain we analyse.

\begin{definition}[The idealized sampler]\label{def:sampler}
A \emph{tangent-space sampler with principal projection} on $\SO(n)$ is a chain
$X_{t+1}=X_tR_{t+1}$ in which the $R_t$ are independent copies of a random rotation $R$, defined
together with a random skew-symmetric matrix $A$ (the tangent proposal, in the frame of the
current point), such that for every outcome either $R=1$ (the proposal is rejected), or
$\skw R=A$ and $\sym R$ is positive semidefinite (the proposal is accepted and projected to the
principal root).
\end{definition}

\begin{remark}[Why this is the sampler of \cite{ZHG}]\label{rem:model}
Consider the sampler of \cite{ZHG} on $\SO(n)$ with uniform target and tangent proposal $V=XA$,
where the entries $A_{ij}$, $i<j$, are independent $N(0,\sigma^2)$; this is the isotropic
Gaussian of scale $s=\sqrt2\,\sigma$ on the tangent space.  Suppose the projection step returns
the principal root whenever it exists and reports failure otherwise.  Then the chain is a
tangent-space sampler with principal projection in the sense of Definition~\ref{def:sampler},
with a conjugation-invariant step distribution.  In Section~\ref{sec:construct} this is proved:
the principal root exists exactly when $1+A^2\succeq0$ and is a measurable function $R(A)$ of
$A$ (Theorem~\ref{thm:principal-exists}); the move of the sampler, with its reverse check and
Metropolis--Hastings test, is $X\mapsto XR(A)$ (Theorem~\ref{thm:zhg-move}); and the law of
$R(A)$ is conjugation invariant (Theorem~\ref{thm:gauss-sampler}).
\notformal\ What remains outside the formal development is the relation between this chain and
a particular implementation.  The projection in \cite{ZHG} is computed by a Newton-type
iteration with a bounded number of steps.  Theorem~\ref{thm:principal-exists}(c) shows that the
iteration with the Jacobian frozen at the base point converges to the principal root whenever
one exists and diverges otherwise; we have not analysed the full Newton iteration, in which the
Jacobian is updated, nor the effect of stopping after finitely many steps or of floating-point
arithmetic.  An implementation whose solver fails on some proposals with $1+A^2\succeq0$, or
returns a different root, is a different chain; if the failures depend only on the conjugacy
class of $A$, the theorems for conjugation-invariant step distributions still apply.
\end{remark}

\subsection{Matrix entries are eigenfunctions}
\begin{proposition}\label{prop:son-haar}
With the topology and Borel $\sigma$-algebra inherited from $\R^{n\times n}$, $\SO(n)$ is a
compact topological group.  Its normalized Haar measure $m_n$ is a probability measure that is
invariant under left and right translations.
\lean{SO.instIsTopologicalGroup,SO.instCompactSpace,haarSO,haarSO.instIsHaarMeasure,haarSO.instIsProbabilityMeasure,haarSO.instIsMulRightInvariant}
\end{proposition}

So the standing assumptions of Section~\ref{sec:compact} hold for $G=\SO(n)$, $m=m_n$, and any
probability measure $\nu$ on $\SO(n)$.

\begin{definition}\label{def:conj}
A measure $\nu$ on $\SO(n)$ is \emph{conjugation invariant} if it is preserved by
$R\mapsto QRQ\transp$ for every orthogonal matrix $Q$.  For such $\nu$ put
$\lambda(\nu)=\frac1n\int\tr R\,d\nu(R)$.
\lean{conjO,ConjInvariant,stdEig}
\end{definition}

\begin{theorem}[The standard representation]\label{thm:entries}
Let $n\ge1$ and let $\nu$ be a finite conjugation-invariant measure on $\SO(n)$.
\begin{enumerate}
\item[(a)] $\int R\,d\nu(R)=\lambda(\nu)\cdot1$.
\item[(b)] For all $i,j$ the function $X\mapsto X_{ij}$ satisfies
$P_\nu X_{ij}=\lambda(\nu)\,X_{ij}$, and hence $P_\nu^tX_{ij}=\lambda(\nu)^tX_{ij}$ for all
$t\ge0$.
\item[(c)] If $n\ge2$ then $\int X_{ij}\,dm_n(X)=0$ for all $i,j$.
\end{enumerate}
\lean{fourierMat_std_eq,markovOp_entry,iterate_markovOp_entry,integral_entry_haar}
\end{theorem}

\begin{proof}
(a) Let $M=\int R\,d\nu$.  By Proposition~\ref{prop:compact-fourier}(c), applied to the inclusion
$\SO(n)\to\R^{n\times n}$ and the map $R\mapsto QRQ\transp$, we have $M=QMQ\transp$ for every
orthogonal $Q$.  Taking for $Q$ the reflection $\operatorname{diag}(1,\dots,-1,\dots,1)$ with
$-1$ in position $i$ gives $M_{ij}=-M_{ij}$ for $j\ne i$, so $M$ is diagonal.  Taking for $Q$
the permutation matrix of the transposition $(i\,j)$ gives $M_{ii}=M_{jj}$.  So $M=c\cdot1$, and
taking traces gives $c=\lambda(\nu)$.  (b) follows from (a) and
Proposition~\ref{prop:compact-fourier}(a).  (c) Pick $k\ne i$ and let $g$ be the diagonal
rotation with entries $-1$ in positions $i$ and $k$ and $1$ elsewhere.  Then
$(gX)_{ij}=-X_{ij}$, and $m_n$ is invariant under $X\mapsto gX$.
\end{proof}

Part (a) is Schur's lemma for the standard representation, proved by hand.  Only invariance
under the finitely many reflections and transpositions is used.  Part (b) says that if
$f(X)=\tr(BX)$ is any linear observable, then along the stationary chain the autocorrelation of
$f(X_t)$ at lag $t$ is exactly $\lambda(\nu)^t$.

\begin{theorem}[Upper bound on the spectral gap]\label{thm:gap-upper}
Let $n\ge2$ and let $\nu$ be a conjugation-invariant probability measure on $\SO(n)$.  Every
Poincar\'e constant $\gamma$ of the walk with step distribution $\nu$ and stationary measure
$m_n$ satisfies
\[
  \gamma\;\le\;1-\lambda(\nu)=1-\frac{\E[\tr R]}{n}.
\]
\lean{gap_le_of_conjInvariant}
\end{theorem}

\begin{proof}
By Theorem~\ref{thm:entries} the entries $X_{1j}$ are mean-zero eigenfunctions with eigenvalue
$\lambda(\nu)$.  Since $\sum_jX_{1j}^2=1$, at least one of them has positive $L^2(m_n)$ norm.
Apply Proposition~\ref{prop:gap-tools}(a).
\end{proof}

\subsection{How large is \texorpdfstring{$1-\lambda$}{1 - lambda}?}

\begin{proposition}\label{prop:one-minus-lambda}
Let $n\ge1$ and let $\nu$ be a probability measure on $\SO(n)$.
\begin{enumerate}
\item[(a)] $1-\lambda(\nu)=\dfrac{1}{2n}\displaystyle\int\fro{1-R}^2\,d\nu(R)$.
\item[(b)] $1-\lambda(\nu)\ge\dfrac1{2n}\displaystyle\int\fro{\skw R}^2\,d\nu(R)$.
\item[(c)] Let $\nu$ be the law of the step $R$ of a tangent-space sampler with principal
projection and tangent proposal $A$ (Definition~\ref{def:sampler}), with $R$ measurable and
$\E\fro{A}^2<\infty$.  Then $1-\lambda(\nu)\le\dfrac1n\,\E\fro{A}^2$.
\end{enumerate}
\lean{one_sub_stdEig_eq_frob,one_sub_stdEig_ge,one_sub_stdEig_le_of_step}
\end{proposition}

\begin{proof}
Integrate the inequalities of Proposition~\ref{prop:step-size}.  In (c), $n-\tr R=0$ on
rejection and $n-\tr R\le\fro{A}^2$ on acceptance.
\end{proof}

In (b), $\skw R$ equals $A$ when the proposal is accepted and $0$ when it is rejected, so (b)
reads $1-\lambda\ge\E[\fro{A}^2;\text{accepted}]/(2n)$.  Thus $1-\lambda(\nu)$ is the mean
squared displacement of one step in the Frobenius norm divided by $2n$, and it lies within a
factor of $2$ of $\E[\fro{A}^2;\text{accepted}]/n$.

\begin{theorem}[Spectral gap of the sampler on $\SO(n)$]\label{thm:sampler}
Let $n\ge2$ and $\sigma\ge0$.  Consider a tangent-space sampler with principal projection on
$\SO(n)$ (Definition~\ref{def:sampler}) such that
\begin{enumerate}
\item[(i)] the step $R$ is measurable and its law $\nu$ is conjugation invariant;
\item[(ii)] the tangent proposal $A$ is skew-symmetric, with $\E[A_{ij}^2]=\sigma^2$ for all
$i\ne j$.
\end{enumerate}
Then $\E\fro{A}^2=n(n-1)\sigma^2$, and every Poincar\'e constant $\gamma$ of the sampler
satisfies
\[
  \gamma\;\le\;1-\lambda(\nu)\;\le\;(n-1)\,\sigma^2 .
\]
\lean{integral_frobSq_of_var,one_sub_stdEig_le_of_step,gap_le_of_conjInvariant,sampler_gap_le}
\end{theorem}

\begin{proof}
Combine Theorem~\ref{thm:gap-upper} with Proposition~\ref{prop:one-minus-lambda}(c).
\end{proof}

By Theorem~\ref{thm:gauss-sampler} below, the hypotheses hold for the Gaussian step
distribution of Definition~\ref{def:gauss-step}, which by Theorem~\ref{thm:zhg-move} is the step
of the sampler of \cite{ZHG} with Gaussian proposals of coordinate variance $\sigma^2$ and
principal projection.  The theorem says that its relaxation time is at
least $1/((n-1)\sigma^2)$ steps.  Independence and Gaussianity of the entries of $A$ are not
used, only their second moments, and the rejection rule is arbitrary.  For small $\sigma$
rejections are rare and Proposition~\ref{prop:one-minus-lambda}(b),(c) pin $1-\lambda(\nu)$
between $\tfrac12(n-1)\sigma^2$ and $(n-1)\sigma^2$, up to the rejection probability; the
simulations of Remark~\ref{rem:numerics} show $1-\lambda(\nu)\approx\tfrac12(n-1)\sigma^2$.  For
$n=2$ this is the $\sigma^2/2$ of Theorem~\ref{thm:so2-size}.

\begin{remark}[Step size and dimension]
\notformal\ For the Gaussian proposal the operator norm of $A$ is about $2\sigma\sqrt n$ when
$n$ is large, so a projection exists with non-negligible probability only if
$\sigma\lesssim1/(2\sqrt n)$.  With $\sigma$ of that order the bound of
Theorem~\ref{thm:sampler} is of order $1$: one step of the idealized sampler moves a constant
fraction of the diameter of $\SO(n)$, in contrast with walks generated by rotations in a single
random plane \cite{Rosenthal}.
\end{remark}

\subsection{The sampler with Gaussian proposals, constructed}\label{sec:construct}
Theorem~\ref{thm:sampler} takes the group step $R$ as given.  We now construct it, and show that
one move of the sampler of \cite{ZHG}, with all its checks, is right multiplication by it.  For a
positive semidefinite matrix $B$ write $\sqrt B$ for its positive semidefinite square root.

\begin{theorem}[Existence and computation of the principal root]\label{thm:principal-exists}
Let $A$ be a real skew-symmetric $n\times n$ matrix.
\begin{enumerate}
\item[(a)] There is an orthogonal matrix $R$ with $\skw R=A$ and $\sym R$ positive semidefinite
if and only if $1+A^2$ is positive semidefinite.  In that case $R$ is unique and
\[
  R=R(A):=\sqrt{1+A^2}+A .
\]
Put $R(A)=1$ when $1+A^2$ is not positive semidefinite or $A$ is not skew-symmetric.  Then for
every matrix $A$, $R(A)\in\SO(n)$ and $\sym R(A)$ is positive semidefinite.
\item[(b)] $R(-A)=R(A)\transp$, and $R(QAQ\transp)=Q\,R(A)\,Q\transp$ for every orthogonal $Q$.
\item[(c)] Let $S_0=0$ and $S_{k+1}=S_k-\tfrac12\bigl((1+A+S_k)\transp(1+A+S_k)-1\bigr)$; this is
Newton's method for the equation $(1+A+S)\transp(1+A+S)=1$ in the symmetric unknown $S$, with
the Jacobian frozen at $S=0$, $A=0$.  Then $1+S_k=P_k$, where $P_0=1$ and
$P_{k+1}=P_k-\tfrac12\bigl(P_k^2-(1+A^2)\bigr)$.  If $1+A^2$ is positive semidefinite, then
$1+A+S_k\to R(A)$ as $k\to\infty$.  Conversely, if the sequence $S_k$ converges, then $1+A^2$ is
positive semidefinite.
\item[(d)] The map $A\mapsto R(A)$ from $\R^{n\times n}$ to itself is Borel measurable.
\end{enumerate}
\lean{exists_principal_iff,principalStep_spec,eq_principalStep,principalStep_mem,posSemidef_symPart_principalStep,principalStep_neg,principalStep_conj,one_add_newtonFrozen,tendsto_visserS,tendsto_visser,tendsto_newtonFrozen,posSemidef_of_tendsto_newtonFrozen,isClosed_principal_domain,measurable_principalStep}
\end{theorem}

\begin{proof}
(c) By induction $P_k=1+S_k$ is a polynomial in $1+A^2$, hence symmetric and commuting with $A$,
and then $(P_k+A)\transp(P_k+A)-1=(P_k-A)(P_k+A)-1=P_k^2-(1+A^2)$, which gives the recursion.
Let $B=1+A^2$ be positive semidefinite.  Since $1-B=A\transp A$ is also positive semidefinite,
$B=U\operatorname{diag}(d)U\transp$ with $U$ orthogonal and $d_i\in[0,1]$, and
$P_k=U\operatorname{diag}(s_k(d_i))U\transp$, where $s_0(b)=1$ and
$s_{k+1}(b)=s_k(b)-\tfrac12(s_k(b)^2-b)$.  For $b\in[0,1]$ one checks by induction that
$\sqrt b\le s_{k+1}(b)\le s_k(b)\le1$, using
$s_{k+1}-\sqrt b=(s_k-\sqrt b)\bigl(1-\tfrac12(s_k+\sqrt b)\bigr)$; so $s_k(b)$ decreases to a limit
$L\ge\sqrt b$ with $L^2=b$, that is, to $\sqrt b$.  Hence $P_k\to\sqrt B$.  Conversely, if
$S_k\to S$ then $L=1+S$ is symmetric and satisfies $L^2=B$, so $B=L\transp L$ is positive
semidefinite.
(a) If $R$ exists, $(\sym R)^2=1+A^2$ by Proposition~\ref{prop:structure}(b), and the square of
a symmetric matrix is positive semidefinite.  Conversely let $P=\sqrt{1+A^2}$.  By (c), $P$ is a
limit of matrices commuting with $A$, so $PA=AP$, and $P+A$ is orthogonal with symmetric part
$P$ and skew part $A$ by Proposition~\ref{prop:structure}(b).  Uniqueness is
Proposition~\ref{prop:principal}(b), and $\det R(A)=1$ is Proposition~\ref{prop:principal}(a).
(b) Both sides have the same skew part and positive semidefinite symmetric part, by
Proposition~\ref{prop:structure}(c) and Proposition~\ref{prop:principal}(c); apply uniqueness.
(d) The set $D$ of skew-symmetric $A$ with $1+A^2\succeq0$ is closed, and $R$ is the pointwise
limit of the functions equal to $P_k(A)+A$ on $D$ and to $1$ elsewhere, each of which is
measurable because $A\mapsto P_k(A)$ is continuous.
\end{proof}

For $n=2$ and $A=aJ$ the condition $1+A^2\succeq0$ reads $|a|\le1$ and
$R(aJ)=R_{\theta(a)}$ (Proposition~\ref{prop:so2-model}).  Since $A$ is normal with purely
imaginary eigenvalues, $1+A^2\succeq0$ says that the operator norm of $A$ is at most $1$.

\begin{theorem}[One move of the sampler]\label{thm:zhg-move}
Let $X$ be an orthogonal and $A$ a skew-symmetric $n\times n$ matrix.  Call an orthogonal matrix
$Y$ a \emph{principal projection of $X+XA$ from $X$} if $Y=X+XA+XS$ for some symmetric $S$ and
$\sym(X\transp Y)$ is positive semidefinite.
\begin{enumerate}
\item[(a)] A principal projection of $X+XA$ from $X$ exists if and only if $1+A^2$ is positive
semidefinite, and then it is unique and equal to $XR(A)$.
\item[(b)] Suppose $1+A^2$ is positive semidefinite and let $Y=XR(A)$.  Then
$\skw(Y\transp X)=-A$, so the tangent proposal of the reverse move is $-YA$; the point $X$ is
the principal projection of $Y-YA$ from $Y$; and $\fro{-YA}=\fro{XA}$.
\item[(c)] Let $g\colon\R\to(0,\infty)$ and define the move $M_g(X,A,u)$ as follows.  If
$X+XA$ has no principal projection from $X$, then $M_g=X$.  Otherwise let $Y$ be the principal
projection and $A'=\skw(Y\transp X)$; if $X$ is the principal projection of $Y+YA'$ from $Y$ and
$u\le g(\fro{YA'}^2)/g(\fro{XA}^2)$, then $M_g=Y$, and otherwise $M_g=X$.  Then for every real
$u\le1$,
\[
  M_g(X,A,u)=X\,R(A).
\]
\end{enumerate}
\lean{IsPrincipalProjection,isPrincipalProjection_iff,isPrincipalProjection_reverse,zhgMove,zhgMove_eq}
\end{theorem}

\begin{proof}
(a) By Proposition~\ref{prop:structure}(a), $Y$ is a principal projection if and only if
$X\transp Y$ is orthogonal with skew part $A$ and positive semidefinite symmetric part; apply
Theorem~\ref{thm:principal-exists}(a).  (b) follows from Proposition~\ref{prop:structure}(c) and
Theorem~\ref{thm:principal-exists}(b): $Y\,R(-A)=XR(A)R(A)\transp=X$.  (c) If there is no
principal projection then $R(A)=1$ by definition.  Otherwise the reverse check succeeds and the
ratio equals $1$ by (b).
\end{proof}

Part (c) is the move of \cite{ZHG} for the uniform target when the proposal density at the
tangent vector $V$ is proportional to $g(\fro{V}^2)$, for instance
$g(r)=e^{-r/(2s^2)}$, and $u$ is uniform on $[0,1]$: project, check that the reverse projection
returns to $X$, and accept with the Metropolis--Hastings probability.  The theorem says that the
two tests never reject, so that the chain is $X_{t+1}=X_tR(A_{t+1})$.

\begin{definition}[The Gaussian step distribution]\label{def:gauss-step}
Let $n\ge0$ and $\sigma\ge0$.  Let $(x_{ij})_{i<j}$ be independent standard Gaussian variables
indexed by the pairs $i<j$ in $\{1,\dots,n\}$, and let $A=\sigma\,\tilde A$, where $\tilde A$ is
the skew-symmetric matrix with $\tilde A_{ji}=x_{ij}$ and $\tilde A_{ij}=-x_{ij}$ for $i<j$.
The \emph{Gaussian step distribution} $\nu_{n,\sigma}$ is the law of $R(A)$, a probability
measure on $\SO(n)$.
\lean{Pairs,skewOf,pairGauss,gaussTangent,gaussStepFun,gaussSOStep}
\end{definition}

\begin{theorem}[The Gaussian sampler]\label{thm:gauss-sampler}
Let $n\ge0$, $\sigma\ge0$, and let $A$ and $\nu_{n,\sigma}$ be as in
Definition~\ref{def:gauss-step}.
\begin{enumerate}
\item[(a)] $R(A)$ is a measurable function of $(x_{ij})$.  For every outcome, either
$R(A)=1$, or $\skw R(A)=A$ and $\sym R(A)$ is positive semidefinite.  Moreover
$\E[A_{ij}^2]=\sigma^2$ for $i\ne j$.  So the chain $X_{t+1}=X_tR(A_{t+1})$ is a tangent-space
sampler with principal projection in the sense of Definition~\ref{def:sampler}.
\item[(b)] The law of $(x_{ij})$ is invariant under every linear map that preserves
$\sum_{i<j}x_{ij}^2$, and $\nu_{n,\sigma}$ is conjugation invariant.
\item[(c)] If $n\ge1$, then $P X_{ij}=\lambda(\nu_{n,\sigma})X_{ij}$ for all $i,j$.  If
$n\ge2$, every Poincar\'e constant $\gamma$ of the walk with step distribution $\nu_{n,\sigma}$
and stationary measure $m_n$ satisfies $\gamma\le1-\lambda(\nu_{n,\sigma})$ and
$\gamma\le(n-1)\sigma^2$.
\item[(d)] For $n=2$, $\nu_{2,\sigma}$ is the law of $R_{\theta(a)}$ with $a\sim N(0,\sigma^2)$.
Hence, for $\sigma>0$, the spectral gap of the Gaussian sampler on $\SO(2)$ is exactly
$1-\lambda_1(\sigma)$ and $\lambda_1(\sigma)$ is its least nonnegative $L^2$ contraction bound.
\end{enumerate}
\lean{measurable_gaussStepFun,gaussStepFun_spec,integral_gaussTangent_sq,pi_gaussian_map_of_isometry,gaussSOStep_conjInvariant,gaussSO_markovOp_entry,gaussSO_gap_le_one_sub_stdEig,gaussSO_gap_le,principalStep_skew2,gaussSOStep_two,gaussSO_two_isSpectralGap,gaussSO_two_isLeast_contraction}
\end{theorem}

\begin{proof}
(a) is Theorem~\ref{thm:principal-exists}(a),(d).  (b) The standard Gaussian measure on a
Euclidean space is invariant under linear isometries.  For $Q$ orthogonal, the map
$\tilde A\mapsto Q\tilde AQ\transp$ is linear on skew-symmetric matrices and preserves
$\fro{\tilde A}^2=2\sum_{i<j}x_{ij}^2$, so $QAQ\transp$ has the same law as $A$, and
$R(QAQ\transp)=QR(A)Q\transp$ by Theorem~\ref{thm:principal-exists}(b).  (c) By (a) and (b),
Theorems~\ref{thm:entries}, \ref{thm:gap-upper} and~\ref{thm:sampler} apply.  (d) For $n=2$,
$A=\sigma x_{12}J$ and $R(aJ)=R_{\theta(a)}$; apply
Theorem~\ref{thm:so2-matrix}(c).
\end{proof}

\section{\texorpdfstring{$\SO(3)$}{SO(3)}: the bound is sharp}\label{sec:sothree}

A rotation in $\SO(3)$ is conjugate to a rotation by an angle $\theta\in[0,\pi]$ about an axis,
and $\tr R=1+2\cos\theta$.  For an integer $l\ge0$ define the \emph{character ratio}
\[
  \kappa_l(\theta)=\frac{\sin\bigl((2l+1)\theta/2\bigr)}{(2l+1)\,\sin(\theta/2)}\quad(\sin(\theta/2)\ne0),
  \qquad\kappa_l(\theta)=1\quad(\sin(\theta/2)=0).
\]
Classically, $(2l+1)\kappa_l(\theta)$ is the character of the irreducible representation of
$\SO(3)$ of dimension $2l+1$ at a rotation of angle $\theta$; for $l=1$ this is the standard
representation and $\kappa_1(\theta)=(1+2\cos\theta)/3=\tr R/3$.

\begin{lemma}\label{lem:sin3}
For every integer $m\ge5$ and every $x\in[0,\pi/4]$, $|\sin(mx)|/m\le\sin(3x)/3$.
\lean{abs_sin_mul_div_le}
\end{lemma}

\begin{proof}
If $x\le\pi/6$ and $mx\le\pi$, concavity of $\sin$ on $[0,\pi]$ gives
$\sin(3x)\ge\tfrac3m\sin(mx)$.  If $x\le\pi/6$ and $mx>\pi$, then
$|\sin(mx)|/m\le1/m<x/\pi\le\tfrac13\cdot\tfrac2\pi\cdot3x\le\sin(3x)/3$ by Jordan's inequality.
If $\pi/6<x\le\pi/4$, then $\sin(3x)\ge\sin(3\pi/4)=\sqrt2/2$ and
$1/m\le1/5<\sqrt2/6$.
\end{proof}

\begin{theorem}[The first character dominates]\label{thm:so3-char}
For every integer $l\ge1$ and every $\theta\in[0,\pi/2]$,
\[
  |\kappa_l(\theta)|\;\le\;\kappa_1(\theta)=\frac{1+2\cos\theta}{3}.
\]
Moreover $|\kappa_l(\theta)|\le1$ for all $l\ge0$ and all real $\theta$.
\lean{abs_so3CharRatio_le,so3CharRatio_one_eq,abs_so3CharRatio_le_one}
\end{theorem}

\begin{proof}
Apply Lemma~\ref{lem:sin3} with $m=2l+1$ and $x=\theta/2$ for $l\ge2$, and divide by
$\sin(\theta/2)>0$.  The last claim is Lemma~\ref{lem:sin}.
\end{proof}

Some restriction on $\theta$ is necessary: $\kappa_1(\pi)=-\tfrac13$ while
$\kappa_2(\pi)=\tfrac15$.  The interval $[0,\pi/2]$, on which the symmetric part of the rotation
is positive semidefinite, is not the largest one on which the inequality holds.

\begin{corollary}\label{cor:so3}
Let $\Theta$ be a random angle with $\Theta\in[0,\pi/2]$ almost surely.  Then for every $l\ge1$,
\[
  \bigl|\E[\kappa_l(\Theta)]\bigr|\;\le\;\E[\kappa_1(\Theta)].
\]
\lean{so3_eigenvalue_dominated}
\end{corollary}

\begin{remark}[The classical picture]\label{rem:so3-classical}
\notformal\ By the Peter--Weyl theorem and the classification of the irreducible representations
of $\SO(3)$, $L^2(\SO(3))$ is the orthogonal direct sum of the spaces $V_l$ of matrix
coefficients of the representation of dimension $2l+1$, $l\ge0$.  A conjugation-invariant walk
acts on $V_l$ as the scalar $\lambda_l=\E[\kappa_l(\Theta)]$, where $\Theta$ is the rotation
angle of the step, by Schur's lemma, and $\lambda_1=\lambda(\nu)$ in the notation of
Section~\ref{sec:son}.  If $\Theta\in[0,\pi/2]$ almost surely, Corollary~\ref{cor:so3} gives
$|\lambda_l|\le\lambda_1$ for all $l\ge1$, so that the spectral gap is $1-\lambda(\nu)$.  None of
these classical inputs is available in Mathlib.  The rest of this section proves the conclusion
(Theorem~\ref{thm:so3-gap}) without them, by constructing on $\mathrm{SU}(2)$ exactly the part of
the theory that is needed; this remark is not used.  In the classical language, the space
$\mathcal H_k$ below is the space of matrix coefficients of the representation
$\operatorname{Sym}^k\C^2$ of $\mathrm{SU}(2)$, and $V_l$ corresponds to $\mathcal H_{2l}$.  One
further consequence that we do not formalize: for Gaussian proposals,
$\cos\Theta=\sqrt{1-|a|^2}$ on acceptance, where $|a|^2=\tfrac12\fro{A}^2$ has mean $3\sigma^2$,
so the gap is $\sigma^2$ to leading order.
\end{remark}

\subsection{Harmonic analysis on \texorpdfstring{$\mathrm{SU}(2)$}{SU(2)} by hand}\label{sec:su2}
Let $\mathrm{SU}(2)$ be the group of complex $2\times2$ matrices $q$ with $q^*q=1$ and
$\det q=1$; it is a compact group, and we write $m$ for its normalized Haar measure.  Every
$q\in\mathrm{SU}(2)$ has the form
$\begin{psmallmatrix}a&b\\-\bar b&\bar a\end{psmallmatrix}$ with $|a|^2+|b|^2=1$.  For
$u,v\in\C^2$ let $\ell_{u,v}(q)=u\transp q\,v$, a linear form in the entries of $q$, and for
$k\ge0$ let
\[
  \mathcal H_k=\operatorname{span}_\C\bigl\{\ell_{u,v}^{\,k}:u,v\in\C^2\bigr\},\qquad
  \mathcal H=\sum_{k\ge0}\mathcal H_k ,
\]
spaces of continuous functions on $\mathrm{SU}(2)$.  A measure $\nu$ on a group $G$ is
\emph{invariant under inner automorphisms} if it is preserved by $q\mapsto gqg^{-1}$ for every
$g\in G$, and for $q_0\in G$ the \emph{class measure} $\nu_{q_0}$ is the image of normalized
Haar measure under $g\mapsto gq_0g^{-1}$.

\begin{proposition}[The spaces $\mathcal H_k$]\label{prop:su2-complete}
The following hold.
\begin{enumerate}
\item[(a)] (Polarization.)  Let $V$ be a commutative $\C$-algebra, $W\subseteq V$ a linear
subspace and $p$ a polynomial with coefficients in $V$.  If $p(\tau)\in W$ for every
$\tau\in\C$, then every coefficient of $p$ lies in $W$.
\item[(b)] For all $u,v,u',v'\in\C^2$ and $k\ge0$,
$\ell_{u',v'}\,\ell_{u,v}^{\,k+1}\in\mathcal H_{k+2}+\mathcal H_k$.
\item[(c)] $\mathcal H$ contains the constants and the matrix entries, and it is closed under
multiplication and under complex conjugation.  Hence it contains every polynomial in the entries
of $q$ and their complex conjugates.
\item[(d)] Every continuous function on $\mathrm{SU}(2)$ is a uniform limit of elements of
$\mathcal H$.
\item[(e)] Every $F\in\mathcal H_k$ is continuous and satisfies $F(-q)=(-1)^kF(q)$.
\end{enumerate}
\lean{coeff_mem_of_forall_eval_mem,harm,lin_mul_lin_pow_mem,polar_two,lin_mul_lin_sub,lin_mul_lin_pow_succ_mem,harmAll,mul_mem_harmAll,conj_mem_harmAll,mem_harmAll_of_mem_polyAlg,polyAlg_topologicalClosure,exists_harmAll_near,continuous_of_mem_harm,harm_parity}
\end{proposition}

\begin{proof}
(a) Otherwise some linear functional vanishes on $W$ but not on a coefficient of $p$; applying
it to $p$ gives a nonzero complex polynomial that vanishes identically.  (b) The function
$\tau\mapsto\ell_{u+\tau u',v}^{\,k+1}=(\ell_{u,v}+\tau\ell_{u',v})^{k+1}$ takes values in
$\mathcal H_{k+1}$, so by (a) its linear coefficient $(k+1)\,\ell_{u',v}\ell_{u,v}^{\,k}$ lies in
$\mathcal H_{k+1}$.  Applying this with $v+\tau v'$ in place of $v$ and using (a) again,
\[
  (k+1)\,\ell_{u,v}^{\,k}\,\ell_{u,v'}\,\ell_{u',v}+\ell_{u,v}^{\,k+1}\,\ell_{u',v'}\in\mathcal H_{k+2}.
\]
Since $\det q=1$, the Cauchy--Binet identity gives
$\ell_{u',v'}\ell_{u,v}-\ell_{u,v'}\ell_{u',v}=(u_1u'_2-u_2u'_1)(v_1v'_2-v_2v'_1)$, a constant.
Substituting, $(k+2)\,\ell_{u',v'}\ell_{u,v}^{\,k+1}$ is the displayed element of
$\mathcal H_{k+2}$ plus a multiple of $\ell_{u,v}^{\,k}$.  (c) By (b) and induction, $\mathcal H$
is stable under multiplication by linear forms, hence by their powers, hence by $\mathcal H$.
The conjugate of $\ell_{u,v}$ is again a linear form because $q_{22}=\overline{q_{11}}$ and
$q_{21}=-\overline{q_{12}}$.  (d) By (c) and the Stone--Weierstrass theorem, since the matrix
entries separate points.  (e) is clear on the generators.
\end{proof}

\begin{theorem}[Schur's lemma for the spaces $\mathcal H_k$, by hand]\label{thm:su2-schur}
Let $\nu$ be a probability measure on $\mathrm{SU}(2)$ invariant under inner automorphisms, and
put $\lambda_k(\nu)=\int q_{11}^{\,k}\,d\nu(q)$.
\begin{enumerate}
\item[(a)] For all $w,v\in\C^2$ and $k\ge0$:
$\displaystyle\int(w\transp q\,v)^k\,d\nu(q)=\lambda_k(\nu)\,(w\transp v)^k$.
\item[(b)] $P_\nu F=\lambda_k(\nu)\,F$ for every $F\in\mathcal H_k$.
\end{enumerate}
\lean{su2Eig,integral_mixed_eq_zero,integral_first_column_pow,exists_su2_first_column,integral_lin_pow,markovOp_lin_pow,markovOp_of_mem_harm}
\end{theorem}

\begin{proof}
(a) We may assume $v\ne0$.  Then $v=r\,ge_1$ for some $g\in\mathrm{SU}(2)$ and $r\ne0$.
Replacing $q$ by $gqg^{-1}$, which does not change the integral, the integrand becomes
$r^k(w'_1q_{11}+w'_2q_{21})^k$ with $w'=g\transp w$.  Expand by the binomial theorem.
Conjugation by $\operatorname{diag}(e^{i\varphi},e^{-i\varphi})$ fixes $q_{11}$ and multiplies
$q_{21}$ by $e^{-2i\varphi}$; taking $\varphi=\pi/(2j)$ shows
$\int q_{11}^{\,a}q_{21}^{\,j}\,d\nu=0$ for $j\ge1$.  What remains is
$(rw'_1)^k\lambda_k(\nu)$, and $rw'_1=w\transp v$.  (b) For $F=\ell_{u,v}^{\,k}$,
$F(xq)=\ell_{x\transp u,v}(q)^k$, so (a) gives $P_\nu F(x)=\lambda_k(\nu)(u\transp xv)^k$.
\end{proof}

\begin{theorem}[The eigenvalues]\label{thm:su2-weyl}
Let $\nu$ be a probability measure on $\mathrm{SU}(2)$ invariant under inner automorphisms.
\begin{enumerate}
\item[(a)] For every $k\ge0$:
$(k+2)\displaystyle\int\tr(q)\,q_{11}^{\,k+1}\,d\nu(q)=(k+3)\,\lambda_{k+2}(\nu)+(k+1)\,\lambda_k(\nu)$.
\item[(b)] $\lambda_0(\nu)=1$ and $2\lambda_1(\nu)=\int\tr q\,d\nu(q)$.  If $\tr q=c$ for
$\nu$-almost every $q$, then
$(k+3)\lambda_{k+2}=(k+2)\,c\,\lambda_{k+1}-(k+1)\lambda_k$ for all $k\ge0$.
\item[(c)] If $\tr q=2\cos\alpha$ for $\nu$-almost every $q$, then
$(k+1)\,\lambda_k(\nu)\,\sin\alpha=\sin\bigl((k+1)\alpha\bigr)$ for all $k\ge0$.
\item[(d)] Let $q_0\in\mathrm{SU}(2)$ with $c=\re(q_0)_{11}$ satisfying $c\ge0$ and $2c^2\ge1$.
Then $|\lambda_{2l}(\nu_{q_0})|\le(4c^2-1)/3$ for every $l\ge1$.
\end{enumerate}
\lean{integral_of_mem_harm,su2Eig_trace_identity,su2Eig_zero,two_mul_su2Eig_one,su2Eig_recurrence,chebS,su2Eig_mul_sin,classMeasure,classMeasure_isConjInvariant,ae_classMeasure_trace,norm_su2Eig_le_one,norm_su2Eig_class_le}
\end{theorem}

\begin{proof}
(a) By the proof of Proposition~\ref{prop:su2-complete}(b) with $u=v=e_1$, $u'=v'=e_2$, the
function $G=(k+1)\,q_{11}^{\,k}q_{12}q_{21}+q_{11}^{\,k+1}q_{22}$ lies in $\mathcal H_{k+2}$.
Evaluating Theorem~\ref{thm:su2-schur}(b) at the identity, $\int G\,d\nu=\lambda_{k+2}G(1)
=\lambda_{k+2}$.  Now use $q_{12}q_{21}=q_{11}q_{22}-1$ and $\tr q=q_{11}+q_{22}$.  (b)
$\lambda_0=1$ is clear; conjugation by $\begin{psmallmatrix}0&1\\-1&0\end{psmallmatrix}$
exchanges $q_{11}$ and $q_{22}$, so $2\lambda_1=\int\tr q\,d\nu$; the recurrence is (a).  (c)
Induction on $k$, using $\sin((k+3)\alpha)=2\cos\alpha\,\sin((k+2)\alpha)-\sin((k+1)\alpha)$.
(d) The trace is constant on a conjugacy class and $\tr q_0=2c$.  Let
$\alpha=\arccos c\in[0,\pi/4]$.  If $\sin\alpha\ne0$, (c) says
$\lambda_{2l}(\nu_{q_0})=\kappa_l(2\alpha)$; if $\sin\alpha=0$ then $\kappa_l(2\alpha)=1$ and
$|\lambda_{2l}|\le1$ because $|q_{11}|\le1$.  In both cases
$|\lambda_{2l}(\nu_{q_0})|\le|\kappa_l(2\alpha)|\le\kappa_1(2\alpha)=(1+2\cos2\alpha)/3
=(4c^2-1)/3$ by Theorem~\ref{thm:so3-char}.
\end{proof}

Part (c) is the Weyl character formula for $\mathrm{SU}(2)$: the eigenvalue of the class measure
of an element with eigenvalues $e^{\pm i\alpha}$ on $\mathcal H_k$ is
$\sin((k+1)\alpha)/((k+1)\sin\alpha)$.

\begin{proposition}[Orthogonality and the $L^2$ bound]\label{prop:su2-l2}
With $m$ the normalized Haar measure of $\mathrm{SU}(2)$, the following hold.
\begin{enumerate}
\item[(a)] If $K\ne j$, $F\in\mathcal H_K$ and $H\in\mathcal H_j$, then
$\int F\,\overline H\,dm=0$.  In particular $\int F\,dm=0$ for $F\in\mathcal H_k$, $k\ge1$.
\item[(b)] Let $\nu$ be a probability measure on $\mathrm{SU}(2)$ invariant under inner
automorphisms, $S$ a finite set of integers, $F_k\in\mathcal H_k$ for $k\in S$, and suppose
$|\lambda_k(\nu)|\le\rho$ for all $k\in S$.  Then $F=\sum_{k\in S}F_k$ satisfies
$\int|P_\nu F|^2\,dm\le\rho^2\int|F|^2\,dm$.
\item[(c)] Let $F$ be a continuous function on $\mathrm{SU}(2)$ with $F(-q)=F(q)$ for all $q$
and $\int F\,dm=0$.  Then $F$ is a uniform limit of finite sums of elements of the spaces
$\mathcal H_{2l}$ with $l\ge1$.
\item[(d)] Let $\nu$ be as in (b), and suppose $|\lambda_{2l}(\nu)|\le\rho$ for all $l\ge1$.
Then $\int|P_\nu F|^2\,dm\le\rho^2\int|F|^2\,dm$ for every $F$ as in (c).  In particular this
holds for $\nu=\nu_{q_0}$ and $\rho=(4c^2-1)/3$ when $c=\re(q_0)_{11}\ge0$ and $2c^2\ge1$.
\end{enumerate}
\lean{classMeasure_map_inv,integral_mul_conj_eq_zero_of_mem_harm,integral_eq_zero_of_mem_harm,integral_norm_sq_sum_of_orthogonal,markovOp_sum_harm,integral_norm_sq_markovOp_sum_harm,exists_const_of_mem_harm_zero,exists_even_harm_near,integral_norm_sq_markovOp_le_of_approx,su2_l2Bound_of_even,su2_class_l2Bound}
\end{proposition}

\begin{proof}
(a) Let $K>j$, $\alpha=\pi/(K+1)$ and $t=\operatorname{diag}(e^{i\alpha},e^{-i\alpha})$.  Since
$t^{-1}$ is conjugate to $t$, the class measure $\nu_t$ is invariant under inversion, so
$P_{\nu_t}$ is symmetric on $L^2(m)$ (Proposition~\ref{prop:compact-basic}).  By
Theorem~\ref{thm:su2-weyl}(c), $\lambda_K(\nu_t)=0$ because $\sin((K+1)\alpha)=0$, and
$\lambda_j(\nu_t)\ne0$ because $0<(j+1)\alpha<\pi$.  As $\overline H\in\mathcal H_j$,
\[
  0=\int(P_{\nu_t}F)\,\overline H\,dm=\int F\,P_{\nu_t}\overline H\,dm
  =\lambda_j(\nu_t)\int F\,\overline H\,dm .
\]
(b) By Theorem~\ref{thm:su2-schur}(b), $P_\nu F=\sum_k\lambda_k(\nu)F_k$, and by (a) and
Pythagoras
\[
  \int|P_\nu F|^2\,dm=\sum_k|\lambda_k(\nu)|^2\int|F_k|^2\,dm\le\rho^2\sum_k\int|F_k|^2\,dm
  =\rho^2\int|F|^2\,dm .
\]
(c) By
Proposition~\ref{prop:su2-complete}(d) there is $h=\sum_kh_k\in\mathcal H$ with
$|F-h|\le\varepsilon$ everywhere.  Since $F$ is even, the even part
$\tfrac12(h(q)+h(-q))=\sum_{k\text{ even}}h_k(q)$ is also within $\varepsilon$ of $F$.  The term
$h_0$ is a constant, equal to the mean of the even part by (a), hence of modulus at most
$\varepsilon$ because $\int F\,dm=0$.  So $\sum_{k\text{ even},\,k\ge2}h_k$ is within
$2\varepsilon$ of $F$.  (d) follows from (b), (c) and the following elementary fact: if
$f\in L^2(m)$ can be approximated in $L^2(m)$ by functions $p$ with
$\int|P_\nu p|^2\le\rho^2\int|p|^2$, then $\int|P_\nu f|^2\le\rho^2\int|f|^2$, because $P_\nu$
does not increase the $L^2$ norm.  The last assertion is Theorem~\ref{thm:su2-weyl}(d).
\end{proof}

\subsection{The double cover}\label{sec:cover}
For real $w,x,y,z$ let
\[
  \varrho(w,x,y,z)=\begin{pmatrix}
  w^2+x^2-y^2-z^2&2(xy-wz)&2(xz+wy)\\
  2(xy+wz)&w^2-x^2+y^2-z^2&2(yz-wx)\\
  2(xz-wy)&2(yz+wx)&w^2-x^2-y^2+z^2\end{pmatrix},
\]
the matrix of $v\mapsto pv\bar p$ on imaginary quaternions for $p=w+xi+yj+zk$.  For
$q\in\mathrm{SU}(2)$ with $q_{11}=w+ix$ and $q_{12}=y+iz$ put $\pi(q)=\varrho(w,x,y,z)$.

\begin{proposition}[The double cover $\mathrm{SU}(2)\to\SO(3)$]\label{prop:cover}
The map $\pi$ has the following properties.
\begin{enumerate}
\item[(a)] $\pi$ maps $\mathrm{SU}(2)$ onto $\SO(3)$; it is a continuous group homomorphism,
and $\pi(-q)=\pi(q)$.
\item[(b)] The image under $\pi$ of the normalized Haar measure of $\mathrm{SU}(2)$ is the
normalized Haar measure $m_3$ of $\SO(3)$.  For every $q_0$,
$\pi_*\nu_{q_0}=\nu_{\pi(q_0)}$.
\item[(c)] $\tr\pi(q)=4(\re q_{11})^2-1$, and if $\sym\pi(q)$ is positive semidefinite then
$2(\re q_{11})^2\ge1$.
\end{enumerate}
\lean{quatRot,quatRot_mul,quatRot_mul_transpose,det_quatRot,covMat_mul,covMat_mem,cover,continuous_cover,cover_negOne_mul,covMat_torus,covMat_yEl,eq_rotX_of_entry,cover_surjective,measurePreserving_cover,classMeasure_cover,trace_covMat,one_le_of_posSemidef_symPart}
\end{proposition}

\begin{proof}
(a) Multiplicativity is a polynomial identity for $\varrho$, as are
$\varrho\varrho\transp=N^2\cdot1$ and $\det\varrho=N^3$ with $N=w^2+x^2+y^2+z^2$; on
$\mathrm{SU}(2)$, $N=1$.  For surjectivity, $\pi$ maps
$\operatorname{diag}(e^{i\varphi},e^{-i\varphi})$ to the rotation by $2\varphi$ about the first
axis and $\begin{psmallmatrix}\cos\beta&\sin\beta\\-\sin\beta&\cos\beta\end{psmallmatrix}$ to the
rotation by $2\beta$ about the second axis.  Given $R\in\SO(3)$, choose a rotation $R_1$ about
the first axis and then a rotation $R_2$ about the second axis such that $M=R_2R_1R$ has
$M_{11}=1$.  Then $M$ fixes $e_1$, so it is a rotation about the first axis, and
$R=R_1^{-1}R_2^{-1}M$ is in the image.  (b) A continuous surjective homomorphism between compact
groups maps normalized Haar measure to normalized Haar measure; the second claim follows because
$\pi(gq_0g^{-1})=\pi(g)\pi(q_0)\pi(g)^{-1}$.  (c) With $N=1$ the trace is $3w^2-x^2-y^2-z^2
=4w^2-1$, and $\sym\pi(q)=(2w^2-1)\cdot1+2nn\transp$ with $n=(x,y,z)$; evaluate the quadratic
form on a nonzero vector orthogonal to $n$.
\end{proof}

\subsection{The spectral gap on \texorpdfstring{$\SO(3)$}{SO(3)}}\label{sec:so3-gap}

\begin{proposition}[From conjugacy classes to invariant measures]\label{prop:classes}
Let $G$ be a compact, second countable topological group with its Borel $\sigma$-algebra, $m$ a
right-invariant probability measure on $G$, and $\nu$ a probability measure on $G$ invariant
under inner automorphisms.  For $R\in G$ let $\nu_R$ be the image of $m$ under
$g\mapsto gRg^{-1}$.
\begin{enumerate}
\item[(a)] (Weighted Cauchy--Schwarz.)  For a measure $\mu$, an integrable complex function $h$
and an integrable function $w>0$ such that $|h|^2/w$ is integrable,
$\bigl|\int h\,d\mu\bigr|^2\le\int w\,d\mu\cdot\int|h|^2/w\,d\mu$.
\item[(b)] Let $f\colon G\to\C$ and $w\colon G\to(0,\infty)$ be continuous, and suppose that
$\int|P_{\nu_R}f|^2\,dm\le w(R)^2\int|f|^2\,dm$ for $\nu$-almost every $R$.  Then
\[
  \int|P_\nu f|^2\,dm\;\le\;\Bigl(\int w\,d\nu\Bigr)^2\int|f|^2\,dm .
\]
\end{enumerate}
\lean{norm_integral_sq_le_weighted,integral_classMeasure,integral_norm_sq_markovOp_le_of_class}
\end{proposition}

\begin{proof}
(a) For every real $s$, $0\le\int(sw-|h|)^2/w\,d\mu=s^2\int w-2s\int|h|+\int|h|^2/w$; so the
discriminant of this quadratic is nonpositive.  (b) By invariance of $\nu$ and Fubini's
theorem, for every $x$,
\[
  P_\nu f(x)=\int\!\!\int f(x\,gRg^{-1})\,d\nu(R)\,dm(g)=\int P_{\nu_R}f(x)\,d\nu(R).
\]
By (a), $|P_\nu f(x)|^2\le\int w\,d\nu\cdot\int|P_{\nu_R}f(x)|^2/w(R)\,d\nu(R)$.  Integrate over
$x$, exchange the integrals, and use the hypothesis.  All functions involved are continuous on
compact spaces, so the integrals exist.
\end{proof}

\begin{theorem}[The spectral gap on $\SO(3)$]\label{thm:so3-gap}
Let $\nu$ be a conjugation-invariant probability measure on $\SO(3)$ such that $\sym R$ is
positive semidefinite for $\nu$-almost every $R$.  Let $\lambda(\nu)=\E[\tr R]/3$.
\begin{enumerate}
\item[(a)] If $R_0\in\SO(3)$ has positive semidefinite symmetric part, then $\tr R_0\ge1$, and
for every continuous $f$ on $\SO(3)$ with $\int f\,dm_3=0$,
$\int|P_{\nu_{R_0}}f|^2\,dm_3\le(\tr R_0/3)^2\int|f|^2\,dm_3$.
\item[(b)] $\lambda(\nu)\ge\tfrac13$, and for every complex $f\in L^2(m_3)$ with
$\int f\,dm_3=0$,
\[
  \int|P_\nu f|^2\,dm_3\;\le\;\lambda(\nu)^2\int|f|^2\,dm_3 ;
\]
moreover $\lambda(\nu)$ is the least nonnegative constant with this property.
\item[(c)] The walk satisfies the Poincar\'e inequality with constant $1-\lambda(\nu)$, and
this is its spectral gap:
\[
  \gap=1-\lambda(\nu)=1-\frac{\E[\tr R]}{3}.
\]
\end{enumerate}
\lean{one_le_trace_of_posSemidef,markovOp_cover,integral_comp_cover,so3_class_l2Bound,ConjInvariant.isConjInvariant,so3_l2Bound_continuous,one_third_le_stdEig,so3_l2ContractionBound,so3_isLeast_contraction,so3_poincareIneq,so3_isSpectralGap}
\end{theorem}

\begin{proof}
(a) Choose $q_0$ with $\pi(q_0)=R_0$ and, replacing $q_0$ by $-q_0$ if necessary,
$c=\re(q_0)_{11}\ge0$.  By Proposition~\ref{prop:cover}(c), $2c^2\ge1$ and
$\tr R_0=4c^2-1\ge1$.  The function $F=f\circ\pi$ is continuous, even, and has mean zero by
Proposition~\ref{prop:cover}(b), and $(P_{\nu_{R_0}}f)\circ\pi=P_{\nu_{q_0}}F$.  So the claim is
Proposition~\ref{prop:su2-l2}(d) with $\rho=(4c^2-1)/3=\tr R_0/3$.  (b) The first claim is (a)
integrated.  For continuous $f$ the inequality follows from (a) and
Proposition~\ref{prop:classes}(b) with $w(R)=\max\{\tr R/3,\tfrac13\}$, which equals $\tr R/3$
for $\nu$-almost every $R$.  A general $f\in L^2(m_3)$ of mean zero is an $L^2$ limit of
continuous functions of mean zero, and the inequality passes to the limit as in the proof of
Proposition~\ref{prop:su2-l2}(d).  Minimality follows from (c).  (c) The Poincar\'e inequality
follows from (b) and Proposition~\ref{prop:gap-tools}(b); the constant is optimal by
Theorem~\ref{thm:gap-upper}.
\end{proof}

\begin{corollary}[The Gaussian sampler on $\SO(3)$]\label{cor:so3-gauss}
For every $n\ge0$ and $\sigma\ge0$, $\sym R$ is positive semidefinite for
$\nu_{n,\sigma}$-almost every $R$.  For every $\sigma\ge0$ the spectral gap of the walk on
$\SO(3)$ with the Gaussian step distribution $\nu_{3,\sigma}$ of Definition~\ref{def:gauss-step}
is exactly
\[
  1-\lambda(\nu_{3,\sigma})=1-\tfrac13\,\E[\tr R(A)],
\]
and $\lambda(\nu_{3,\sigma})$ is its least nonnegative $L^2$ contraction bound.
\lean{ae_posSemidef_gaussSOStep,gaussSO_three_isSpectralGap,gaussSO_three_isLeast_contraction}
\end{corollary}

\begin{proof}
Theorem~\ref{thm:principal-exists}(a), Theorem~\ref{thm:gauss-sampler}(b) and
Theorem~\ref{thm:so3-gap}.
\end{proof}

By Theorem~\ref{thm:gauss-sampler}(c) this gap is at most $2\sigma^2$.  The hypothesis on the
symmetric part in Theorem~\ref{thm:so3-gap} cannot be dropped: for the class measure of a
rotation by $\pi$, $\lambda=-\tfrac13$ while the eigenvalue on $\mathcal H_4$ is
$\kappa_2(\pi)=\tfrac15$.

\section{The spectral gap on \texorpdfstring{$\SO(n)$}{SO(n)}}\label{sec:son-gap}

For the Gaussian sampler on $\SO(2)$ and $\SO(3)$ the spectral gap is $1-\lambda(\nu)$, and on
$\SO(3)$ this is a consequence of conjugation invariance and the positivity of the symmetric
part alone (Theorem~\ref{thm:so3-gap}).  The next two theorems show that on $\SO(4)$ it is not.
The same holds for every even $n$, including $n=2$; for odd $n\ge5$ we do not know (see the
discussion after Theorem~\ref{thm:so4}).  The rest of the section shows that it is a consequence of a different
hypothesis, which the Gaussian sampler satisfies for every $n$: the step distribution has a
density with respect to Haar measure that decreases with each rotation angle
(Theorem~\ref{thm:monotone} and Corollary~\ref{cor:son-gauss}).

\subsection{A second family of eigenfunctions}
For indices $a,b,i,j$ let $\mu^{ab}_{ij}(X)=X_{ai}X_{bj}-X_{aj}X_{bi}$, the $2\times2$ minor of
$X$ with rows $a,b$ and columns $i,j$.

\begin{theorem}[Two-by-two minors]\label{thm:second}
Let $\nu$ be a finite conjugation-invariant measure on $\SO(n)$, let $i\ne j$, and put
\[
  \lambda_2^{ij}(\nu)=\int\bigl(R_{ii}R_{jj}-R_{ij}R_{ji}\bigr)\,d\nu(R).
\]
\begin{enumerate}
\item[(a)] For all $k,l$:
$\int\mu^{kl}_{ij}(R)\,d\nu(R)=\lambda_2^{ij}(\nu)\,(\delta_{ki}\delta_{lj}-\delta_{kj}\delta_{li})$.
\item[(b)] For all $a,b$: $P_\nu\,\mu^{ab}_{ij}=\lambda_2^{ij}(\nu)\,\mu^{ab}_{ij}$.
\item[(c)] $\lambda_2^{ij}(\nu)=\lambda_2^{kl}(\nu)$ for all $i\ne j$ and $k\ne l$.  We write
$\lambda_2(\nu)$ for the common value.
\item[(d)] If $n\ge3$, then $\int\mu^{ab}_{ij}\,dm_n=0$ and $\int(\mu^{ij}_{ij})^2\,dm_n>0$.
Hence, if moreover $\nu$ is a probability measure, every Poincar\'e constant $\gamma$ of the
walk satisfies
\[
  \gamma\;\le\;1-\lambda_2(\nu).
\]
\end{enumerate}
\lean{minor2,minorEig,integral_conjO,integral_minor2,markovOp_minor2,minorEig_comm,minorEig_swap,minorEig_eq,integral_minor2_haar,integral_minor2_sq_pos,gap_le_one_sub_minorEig}
\end{theorem}

\begin{proof}
(a) Conjugation by the reflection in coordinate $c$ multiplies $\mu^{kl}_{ij}(R)$ by the product
of the signs of $k,l,i,j$, where the sign of an index is $-1$ if it equals $c$ and $1$
otherwise.  If $k\ne l$ and $\{k,l\}\ne\{i,j\}$, some $c\in\{k,l\}$ occurs exactly once among
$k,l,i,j$, so the integral equals its negative.  If $k=l$ the minor vanishes, and
$\mu^{ji}_{ij}=-\mu^{ij}_{ij}$.  (b) Expanding the product,
$\mu^{ab}_{ij}(XR)=\sum_{k,l}X_{ak}X_{bl}\,\mu^{kl}_{ij}(R)$; integrate and use (a).  (c)
Conjugation by the permutation matrix of a transposition $\tau$ gives
$\lambda_2^{\tau(i)\tau(j)}=\lambda_2^{ij}$, the definition is symmetric in $i,j$, and two
transpositions carry $(i,j)$ to $(k,l)$.  (d) Choose $c\notin\{i,j\}$ and let
$g$ be the diagonal rotation with entries $-1$ in positions $i$ and $c$; then
$\mu^{ab}_{ij}(Xg)=-\mu^{ab}_{ij}(X)$, and $m_n$ is right invariant.  The function
$(\mu^{ij}_{ij})^2$ is continuous, nonnegative and equal to $1$ at the identity.  Apply
Proposition~\ref{prop:gap-tools}(a).
\end{proof}

Part (a) is Schur's lemma for the second exterior power of the standard representation, proved
by hand as in Theorem~\ref{thm:entries}; as there, only reflections and transpositions are used.
Summing over the pairs, $\lambda_2(\nu)=\E[e_2(R)]/\binom n2$, where
$e_2(R)=\tfrac12\bigl((\tr R)^2-\tr(R^2)\bigr)$ is the second elementary symmetric function of the
eigenvalues.  Together with Theorem~\ref{thm:gap-upper},
$\gap\le1-\max\{\lambda(\nu),\lambda_2(\nu)\}$ for $n\ge3$.

\subsection{A walk on \texorpdfstring{$\SO(4)$}{SO(4)} whose gap is not given by the standard representation}
The orthogonal group $\Or(n)$ is compact; let $m_{\Or(n)}$ be its normalized Haar measure.  For
$R_0\in\SO(n)$ the \emph{orbit measure} of $R_0$ is the image of $m_{\Or(n)}$ under
$Q\mapsto QR_0Q\transp$, the uniform distribution on the $\Or(n)$-conjugacy class of $R_0$.

\begin{theorem}[The theorem for $\SO(3)$ fails on $\SO(4)$]\label{thm:so4}
Every orbit measure is a conjugation-invariant probability measure.  Let
\[
  J=\begin{pmatrix}0&-1&0&0\\1&0&0&0\\0&0&0&-1\\0&0&1&0\end{pmatrix}\in\SO(4),
\]
the rotation by $\pi/2$ in the coordinate planes $\{1,2\}$ and $\{3,4\}$, and let $\nu_J$ be its
orbit measure.
\begin{enumerate}
\item[(a)] $\sym J=0$, and $\sym R$ is positive semidefinite for $\nu_J$-almost every $R$.
\item[(b)] $\lambda(\nu_J)=0$ and $\lambda_2(\nu_J)=\tfrac13$.
\item[(c)] Every Poincar\'e constant of the walk on $\SO(4)$ with step distribution $\nu_J$ is
at most $\tfrac23$.  In particular the Poincar\'e inequality with constant
$1-\lambda(\nu_J)=1$ is false: the spectral gap is strictly smaller than
$1-\E[\tr R]/4$.
\end{enumerate}
\lean{haarO,orbitMeasure,orbitMeasure_conjInvariant,J4,symPart_J4mat,so4Step,so4Step_conjInvariant,ae_posSemidef_so4Step,stdEig_so4Step,e2_eq,e2_conj_J4,minorEig_so4Step,so4_poincare_le,so4_counterexample}
\end{theorem}

\begin{proof}
Invariance of the orbit measure follows from the left invariance of $m_{\Or(n)}$.  (a)
$J\transp=-J$, and $\sym(QJQ\transp)=Q(\sym J)Q\transp$.  (b) $\tr(QJQ\transp)=\tr J=0$.  For a $4\times4$ matrix the sum of the six
principal $2\times2$ minors is $e_2(R)=\tfrac12((\tr R)^2-\tr(R^2))$, which is invariant under
conjugation, and $e_2(J)=\tfrac12(0-\tr(-1))=2$.  By Theorem~\ref{thm:second}(c) the six numbers
$\lambda_2^{ij}(\nu_J)$ are equal, and their sum is $\int e_2\,d\nu_J=2$.  (c)
Theorem~\ref{thm:second}(d).
\end{proof}

The mechanism is visible in the eigenvalues.  A rotation of $\R^4$ by the angle $\alpha$ in two
orthogonal planes has $\tr R/4=\cos\alpha$ and $e_2(R)/6=(1+2\cos^2\alpha)/3$, and the second
number is the larger one as soon as $\cos\alpha<\tfrac12$.  For $\SO(3)$ this cannot happen:
there $e_2(R)=\tr R$, so $\lambda_2=\lambda$.

The example has an analogue in every even dimension.  For $n=2m\ge4$ let $J_n\in\SO(n)$ be the
block-diagonal matrix with $m$ blocks $R_{\pi/2}$.  Then $\sym J_n=0$, $\tr J_n=0$ and
$e_2(J_n)=m$, so the orbit measure of $J_n$ has $\lambda=0$ and
$\lambda_2=m/\binom n2=1/(n-1)$, and its spectral gap is at most $1-1/(n-1)<1$ by
Theorem~\ref{thm:second}(d).  For $n=2$ the minors are constant, but the step distribution
$\tfrac12(\delta_{R_{\pi/2}}+\delta_{R_{-\pi/2}})$, which is conjugation invariant with
$\sym R=0$ and $\lambda=0$, corresponds on $\T$ to $\nu=\tfrac12(\delta_{\pi/2}+\delta_{-\pi/2})$
with $\hat\nu(k)=\cos(k\pi/2)$ and $\hat\nu(4)=1$, so its spectral gap is $0$
(Theorems~\ref{thm:circle}(d) and~\ref{thm:so2-matrix}(b)).  Thus for even $n$ conjugation invariance and the positivity of
the symmetric part do not imply that the spectral gap is $1-\lambda(\nu)$.

For odd $n$ the minors cannot give such an example.  Let $n=2m+1$ and let $R$ have rotation
angles $\theta_k$ with $c_k=\cos\theta_k\in[0,1]$.  Then $\tr R=1+2\sum_kc_k$ and
$e_2(R)=m+2\sum_kc_k+4\sum_{k<l}c_kc_l$, so
\[
  \frac{\tr R}{n}-\frac{e_2(R)}{\binom n2}
  =\frac{2}{mn}\Bigl((m-1)\sum_kc_k-2\sum_{k<l}c_kc_l\Bigr)\;\ge\;0,
\]
because $2\sum_{k<l}c_kc_l=\sum_kc_k\sum_{l\ne k}c_l\le(m-1)\sum_kc_k$.  Hence
$\lambda_2(\nu)\le\lambda(\nu)$ whenever $n$ is odd and $\sym R$ is positive semidefinite almost
surely.  We do not know whether Theorem~\ref{thm:so3-gap} extends to odd $n\ge5$.  Numerically,
for $n\in\{5,7,9\}$ none of the irreducible representations that we tested has a character ratio
larger in absolute value than that of the standard representation at a rotation with positive
semidefinite symmetric part (Remark~\ref{rem:numerics}(9)); if this holds for every
representation, the theorem extends to these $n$ by the argument of
Remark~\ref{rem:so3-classical}.

\subsection{Step distributions that decrease with the rotation angles}\label{sec:monotone}
Let $m=\lfloor n/2\rfloor$, let $R_a$ be the rotation of $\R^2$ by the angle $a$, and for
$\theta\in\R^m$ let $t(\theta)\in\SO(n)$ be the block-diagonal matrix with diagonal blocks
$R_{\theta_1},\dots,R_{\theta_m}$, followed by a $1$ if $n$ is odd.  Every $R\in\SO(n)$ is
$Qt(\theta)Q\transp$ for some $Q\in\Or(n)$ and $\theta\in[0,\pi]^m$; the numbers $\theta_k$, the
\emph{rotation angles} of $R$, are determined by $R$ up to order, because the eigenvalues of $R$
are $e^{\pm i\theta_k}$ (and $1$ if $n$ is odd).  Conjugating $t(\theta)$ by permutation matrices
and by diagonal matrices with entries $\pm1$ permutes the $\theta_k$ and changes their signs.  So a
function $f$ on $\SO(n)$ that is invariant under conjugation by $\Or(n)$ is determined by
$F(\theta)=f(t(\theta))$, which is symmetric in $\theta$, and even and $2\pi$-periodic in each
$\theta_k$.  The symmetric part of $R$ is positive semidefinite exactly when $\cos\theta_k\ge0$
for all $k$.

We call a probability measure $\nu$ on $\SO(n)$ \emph{angle-monotone} if
$\nu=p\,\delta_1+(1-p)\,f\,m_n$ with $0\le p\le1$ and $f$ a Lipschitz probability density on
$\SO(n)$, invariant under conjugation by $\Or(n)$, such that $F(\theta)=f(t(\theta))$ is
nonincreasing on $[0,\pi]^m$ for the coordinatewise order.  For $n=2$ this says that the step
angle has an even density that decreases on $[0,\pi]$, as in Theorem~\ref{thm:unimodal}.  The
orbit measure $\nu_J$ of Theorem~\ref{thm:so4} is not angle-monotone.

The proof that angle-monotone walks have spectral gap $1-\lambda(\nu)$ integrates by parts on the
group, which turns $\int\chi\,d\nu$, for an eigenfunction $\chi$ of the Laplacian with eigenvalue
$-c$, into $c^{-1}\int\ip{\nabla f}{\nabla\chi}$; it then compares gradients on the torus
$\{t(\theta)\}$.  The comparison is the following lemma, in which $H$ will be a matrix coefficient
averaged over conjugation.

\begin{lemma}[Gradient bound]\label{lem:grad}
Let $m\ge1$, let $w_1,\dots,w_N\ge0$ and $\mu^{(1)},\dots,\mu^{(N)}\in\Z^m$, and put
\[
  H(\theta)=\sum_{a=1}^Nw_a\prod_{k=1}^m\cos\bigl(\mu^{(a)}_k\theta_k\bigr),\qquad\theta\in\R^m .
\]
\begin{enumerate}
\item[(a)] For every $\theta\in\R^m$ and every $k$,
$\displaystyle\Bigl|\frac{\partial H}{\partial\theta_k}(\theta)\Bigr|\le
\Bigl(\sum_aw_a\bigl(\mu^{(a)}_k\bigr)^2\Bigr)\,|\sin\theta_k|$.
\item[(b)] Let $n\ge2$, $d=n(n-1)/2$ and $c>0$, and let $p,q\in\R^m$ satisfy
$p_k\sin\theta_k\le0$ and $|q_k|\le(c/d)\,|\sin\theta_k|$ for every $k$.  Then
\[
  \frac1c\,\Bigl|\sum_kp_kq_k\Bigr|\;\le\;\frac1{n-1}\sum_kp_k\Bigl(-\frac2n\sin\theta_k\Bigr).
\]
\end{enumerate}
\lean{abs_sin_int_mul_le,cosProd,cosSum,cosSumDeriv,cosProd_update,hasDerivAt_cosSum_update,abs_cosSumDeriv_le,abs_sum_mul_le_of_sin,gradient_domination}
\end{lemma}

\begin{proof}
(a) The derivative is the sum over $a$ of
$-w_a\mu_k\sin(\mu_k\theta_k)\prod_{l\ne k}\cos(\mu_l\theta_l)$, with $\mu=\mu^{(a)}$.  The product
of cosines has absolute value at most $1$, and
$|\mu_k\sin(\mu_k\theta_k)|\le\mu_k^2\,|\sin\theta_k|$ by Lemma~\ref{lem:sin} applied to
$|\mu_k|$.  (b) $|\sum_kp_kq_k|\le\sum_k|p_k|(c/d)|\sin\theta_k|=(c/d)\sum_k(-p_k\sin\theta_k)$,
and $\frac1{n-1}\cdot\frac2n=\frac1d$.
\end{proof}

For $H(\theta)=\tr t(\theta)/n=(2\sum_k\cos\theta_k+n-2m)/n$, (a) is an equality.

\begin{lemma}[The projection in angle coordinates]\label{lem:jacobian}
\begin{enumerate}
\item[(a)] Let $a,b\in\R$ and $J=R_{\pi/2}$.  Then $\frac12(R_aJ+JR_a\transp)=\cos a\cdot J$; the
linear map $v\mapsto\frac12(R_av+v)$ of $\R^2$ has determinant $\frac12(1+\cos a)$; and the linear
map $M\mapsto\frac12(R_aM+MR_b\transp)$ of $\R^{2\times2}$ has determinant
$\bigl(\frac12(\cos a+\cos b)\bigr)^2$.
\item[(b)] Let $m\ge1$, $v>0$ and $\epsilon\in\{0,1\}$.  For $c\in[0,\infty)^m$ put
\[
  D_v(c)=\exp\Bigl(-\sum_k\frac{1-c_k^2}{2v}\Bigr)\prod_kc_k\prod_{k<l}\Bigl(\frac{c_k+c_l}2\Bigr)^2
  \prod_k\Bigl(\frac{1+c_k}2\Bigr)^{\epsilon},
\]
and for $\theta\in\R^m$ put $F_v(\theta)=D_v(\cos\theta_1,\dots,\cos\theta_m)$ if $\cos\theta_k\ge0$
for all $k$, and $F_v(\theta)=0$ otherwise.  Then $D_v$ is nondecreasing on $[0,\infty)^m$, and
$F_v\ge0$ is nonincreasing on $[0,\pi]^m$, for the coordinatewise orders.
\end{enumerate}
\lean{rotMat,skew2,skew_rotMat_mul_skew2,det_half_rotMat_add_one,pairBlock,pairBlock_mulVec,det_pairBlock,gaussCosDensity,gaussAngleDensity,gaussCosDensity_nonneg,gaussCosDensity_mono,gaussAngleDensity_nonneg,gaussAngleDensity_antitoneOn}
\end{lemma}

\begin{proof}
(a) Direct computation.  In the basis $E_{11},E_{12},E_{21},E_{22}$ the last map is
$\frac12\bigl((\cos a+\cos b)\cdot1+S\bigr)$ with $S$ skew-symmetric, and its determinant is
$\frac1{16}\bigl((\cos a+\cos b)^2+(\sin a+\sin b)^2\bigr)\bigl((\cos a+\cos b)^2+(\sin a-\sin b)^2\bigr)
=\frac1{16}\bigl(2+2\cos(a-b)\bigr)\bigl(2+2\cos(a+b)\bigr)=\frac14(\cos a+\cos b)^2$.
(b) Each factor of $D_v$ is nonnegative and nondecreasing in each $c_k\ge0$.  If
$0\le\theta\le\theta'\le\pi$ coordinatewise, then $\cos\theta'_k\le\cos\theta_k$ for all $k$; if all
$\cos\theta'_k\ge0$, then $F_v(\theta')=D_v(\cos\theta')\le D_v(\cos\theta)=F_v(\theta)$, and
otherwise $F_v(\theta')=0\le F_v(\theta)$.
\end{proof}

\begin{theorem}[Angle-monotone step distributions]\label{thm:monotone}
\notformal\ Let $n\ge2$, let $\nu$ be an angle-monotone probability measure on $\SO(n)$, and let
$\lambda(\nu)=\E[\tr R]/n$.
\begin{enumerate}
\item[(a)] $\lambda(\nu)\ge0$, and for every complex $f\in L^2(m_n)$ with $\int f\,dm_n=0$,
\[
  \int|P_\nu f|^2\,dm_n\;\le\;\lambda(\nu)^2\int|f|^2\,dm_n ;
\]
moreover $\lambda(\nu)$ is the least nonnegative constant with this property.
\item[(b)] The walk satisfies the Poincar\'e inequality with constant $1-\lambda(\nu)$, and this
is its spectral gap:
\[
  \gap=1-\lambda(\nu)=1-\frac{\E[\tr R]}{n}.
\]
\end{enumerate}
\end{theorem}

\begin{proof}
Put $d=n(n-1)/2$ and $\ip XY=-\frac12\tr(XY)$ on the space $\mathfrak{so}(n)$ of skew-symmetric
matrices.  The matrices $J_{ij}=e_je_i\transp-e_ie_j\transp$, $i<j$, form an orthonormal basis,
and with $E_k=J_{2k-1,2k}$ we have $t(\theta)=\exp(\sum_k\theta_kE_k)$.  For a function $u$ on
$\SO(n)$ and $X\in\mathfrak{so}(n)$ let $(Xu)(R)=\frac d{ds}u(Re^{sX})|_{s=0}$, and put
$\ip{\nabla u}{\nabla w}=\sum_{i<j}(J_{ij}u)(J_{ij}w)$ and $\Delta u=\sum_{i<j}J_{ij}(J_{ij}u)$.
Both are unchanged if $(J_{ij})$ is replaced by another orthonormal basis, for instance
$(QJ_{ij}Q\transp)$ with $Q\in\Or(n)$; hence $\Delta$ commutes with conjugation by orthogonal
matrices.

\emph{Step 1: Green's formula.}  Since $m_n$ is right invariant,
$s\mapsto\int u(Re^{sX})w(Re^{sX})\,dm_n(R)$ is constant.  If $u$ is Lipschitz and $w$ is smooth,
the difference quotients of $u$ along $s\mapsto Re^{sX}$ are bounded and, by Rademacher's theorem,
converge almost everywhere, so differentiating at $s=0$ gives $\int(Xu)w\,dm_n=-\int u(Xw)\,dm_n$.
Hence $\int u\,\Delta w\,dm_n=-\int\ip{\nabla u}{\nabla w}\,dm_n$.

\emph{Step 2: reduction to irreducible subspaces.}  By the Peter--Weyl theorem $L^2(m_n)$ is the
orthogonal sum of finite-dimensional spaces $V$ of smooth functions that are invariant under the
right translations $(\rho(Q)g)(x)=g(xQ)$ and on which $\rho$ is irreducible; the constants are the
only such $V$ on which $\rho$ is trivial.  The operator $P_\nu=\int\rho(R)\,d\nu(R)$ maps each $V$
into itself and commutes with $\rho(Q)$ for $Q\in\SO(n)$, because $\nu$ is invariant under
conjugation; by Schur's lemma it acts on $V$ as a scalar $\lambda_V$.  We show below that
$|\lambda_V|\le\lambda(\nu)$ whenever $V$ is not the constants, and that $\lambda(\nu)\ge0$.  Then
for $f=\sum_Vf_V$ of mean zero,
$\int|P_\nu f|^2=\sum_V|\lambda_V|^2\int|f_V|^2\le\lambda(\nu)^2\int|f|^2$, which is (a); no
smaller constant works because the matrix entries, which have mean zero, are eigenfunctions with
eigenvalue $\lambda(\nu)$ (Theorem~\ref{thm:entries}).  Part (b) follows from (a) by
Proposition~\ref{prop:gap-tools}(b) and Theorem~\ref{thm:gap-upper}.

\emph{Step 3: the Casimir operator.}  Fix $V$ other than the constants and write $D_Xg=Xg$ for
$g\in V$; $D_X$ is a skew-adjoint operator on $V$, and $\rho(Q)D_X\rho(Q)^{-1}=D_{QXQ\transp}$ for
$Q\in\SO(n)$.  So $\Omega=-\Delta|_V=-\sum_{i<j}D_{J_{ij}}^2$ commutes with $\rho$, and it is a
scalar $c\ge0$ by Schur's lemma.  If $c=0$ then $D_Xg=0$ for all $X$ and all $g\in V$, so every
$g\in V$ is invariant under right translation by the connected group $\SO(n)$, that is, constant;
hence $c>0$.  For each $k$ and each pair $i<j$ there is $P\in\SO(n)$ with $PE_kP\transp=J_{ij}$,
so by the invariance of $m_n$ the operator $\int\rho(Q)D_X^2\rho(Q)^{-1}\,dm_n(Q)$ is the same for
$X=E_k$ as for $X=J_{ij}$.  Summing over the $d$ pairs, it equals $-\Omega/d$:
\begin{equation}\label{eq:casimir-avg}
  \int\rho(Q)\,D_{E_k}^2\,\rho(Q)^{-1}\,dm_n(Q)=-\frac cd\quad\text{on }V,\qquad k=1,\dots,m.
\end{equation}

\emph{Step 4: the eigenvalue as a pairing of gradients.}  Let $g\in V$ with $\int|g|^2dm_n=1$, let
$h(R)=\int\overline{g(x)}\,g(xR)\,dm_n(x)$, let $r=\operatorname{diag}(1,\dots,1,-1)$, and put
\[
  \tilde h(R)=\frac12\int\bigl(h(QRQ^{-1})+h(QrRr^{-1}Q^{-1})\bigr)\,dm_n(Q).
\]
Then $\lambda_V=\int\overline g\,P_\nu g\,dm_n=\int h\,d\nu=\int\tilde h\,d\nu$, the last step because $\nu$
is invariant under conjugation by $\Or(n)$, which is generated by $\SO(n)$ and $r$.  The function
$\tilde h$ is smooth and invariant under conjugation by $\Or(n)$, and $\Delta\tilde h=-c\,\tilde h$
because $\Delta h(R)=\int\overline g\,\rho(R)(\Delta g)\,dm_n=-c\,h(R)$.  Likewise $\psi(R)=\tr R/n$
satisfies $\Delta\psi=-(n-1)\psi$, since $\sum_{i<j}J_{ij}^2=-(n-1)\cdot1$.  As
$\tilde h(1)=\psi(1)=1$, Step~1 gives
\[
  \lambda_V=p+\frac{1-p}c\int\ip{\nabla f}{\nabla\tilde h}\,dm_n,\qquad
  \lambda(\nu)=p+\frac{1-p}{n-1}\int\ip{\nabla f}{\nabla\psi}\,dm_n .
\]
So it suffices to show that for $m_n$-almost every $R$
\begin{equation}\label{eq:pointwise}
  \frac1c\,\bigl|\ip{\nabla f}{\nabla\tilde h}(R)\bigr|\;\le\;\frac1{n-1}\,\ip{\nabla f}{\nabla\psi}(R).
\end{equation}

\emph{Step 5: gradients of class functions.}  Let $u$ and $w$ be smooth or Lipschitz functions
invariant under conjugation by $\Or(n)$, and $U=u\circ t$, $W=w\circ t$.  Every $R\in\SO(n)$ is
$Pt(\theta)P^{-1}$ with $P\in\SO(n)$ (if $\det Q=-1$ in $R=Qt(\theta)Q\transp$, replace $Q$ by
$Q\operatorname{diag}(1,-1,1,\dots,1)$ and $\theta_1$ by $-\theta_1$).  Since
$(Xu)(R)=((P^{-1}XP)u)(t(\theta))$, $\ip{\nabla u}{\nabla w}(R)$ equals
$\ip{\nabla u}{\nabla w}(t)$ with $t=t(\theta)$.  Suppose $u$ and $w$ are differentiable at $t$.
For $Z\in\mathfrak{so}(n)$, $u(e^{-sZ}te^{sZ})=u(t)$ and $e^{-sZ}te^{sZ}=t\,e^{-st^{-1}Zt}e^{sZ}$, so
$((Z-t^{-1}Zt)u)(t)=0$, and likewise for $w$.  If $\theta$ is \emph{regular}, meaning
$\cos\theta_k\ne\cos\theta_l$ for $k<l$ and, for odd $n$, $\cos\theta_k\ne1$, then
$Z\mapsto Z-t^{-1}Zt$ maps the orthogonal complement of $\operatorname{span}\{E_k\}$ onto itself,
so $(Xu)(t)=(Xw)(t)=0$ for $X$ in that complement.  Since $t(\theta)e^{sE_k}=t(\theta+se_k)$,
$(E_ku)(t)=\partial_kU(\theta)$.  Hence, at regular points,
\[
  \ip{\nabla u}{\nabla w}(R)=\sum_k\partial_kU(\theta)\,\partial_kW(\theta).
\]
The non-regular elements are the zeros of
$\prod_{k<l}(\cos\theta_k-\cos\theta_l)^2\prod_k(1-\cos\theta_k)^{n-2m}$, a symmetric function of
the eigenvalues and hence a polynomial in the entries of $R$, which does not vanish identically;
so they form a null set.  By Rademacher's theorem $f$ is differentiable almost everywhere.

\emph{Step 6: the bound on the torus.}  The operators $\rho(t(\theta))$, $\theta\in\R^m$, are unitary
and commute, and $t(\theta+2\pi e_k)=t(\theta)$, so $V$ is the orthogonal sum of joint eigenspaces
$V_\mu$ on which $\rho(t(\theta))=e^{i\ip\mu\theta}$, for $\mu$ in a finite set $M\subset\Z^m$; on
$V_\mu$, $D_{E_k}=i\mu_k$.  For $Q\in\SO(n)$ let $a_\mu(Q)$ be the squared norm of the component
of $\rho(Q)^{-1}g$ in $V_\mu$.  Then $h(Qt(\theta)Q^{-1})=\sum_\mu a_\mu(Q)e^{i\ip\mu\theta}$ and
$\sum_\mu a_\mu(Q)\mu_k^2=-\int\overline{\rho(Q)^{-1}g}\;D_{E_k}^2\rho(Q)^{-1}g\,dm_n$.  For a sign
vector $\varepsilon\in\{\pm1\}^m$, $t(\varepsilon\theta)$ is conjugate to $t(\theta)$ by a diagonal
orthogonal matrix, and $rt(\theta)r^{-1}$ is such a $t(\varepsilon\theta)$.  Averaging over the
$2^m$ sign vectors therefore gives
\[
  \tilde H(\theta):=\tilde h(t(\theta))=\sum_{\mu\in M}w_\mu\prod_k\cos(\mu_k\theta_k),\qquad
  w_\mu=\int a_\mu(Q)\,dm_n(Q)\ge0,
\]
and $\sum_\mu w_\mu\mu_k^2=c/d$ by \eqref{eq:casimir-avg}.  By Lemma~\ref{lem:grad}(a),
$|\partial_k\tilde H(\theta)|\le(c/d)\,|\sin\theta_k|$.  For $\Psi=\psi\circ t$,
$\partial_k\Psi(\theta)=-\frac2n\sin\theta_k$.  Finally $F=f\circ t$ is even, $2\pi$-periodic and
nonincreasing on $[0,\pi]$ in each $\theta_k$, so $\partial_kF(\theta)\sin\theta_k\le0$ wherever
$F$ is differentiable.  At every regular $R$ at which $f$ is differentiable, Step~5 and
Lemma~\ref{lem:grad}(b) with $p_k=\partial_kF(\theta)$ and $q_k=\partial_k\tilde H(\theta)$ give
\eqref{eq:pointwise}.  The same computation shows $\ip{\nabla f}{\nabla\psi}\ge0$ almost
everywhere, so $\lambda(\nu)\ge p\ge0$.
\end{proof}

\begin{corollary}[The Gaussian sampler on $\SO(n)$]\label{cor:son-gauss}
\notformal\ Let $n\ge2$, $\sigma\ge0$, and let $\nu_{n,\sigma}$ be the Gaussian step distribution
of Definition~\ref{def:gauss-step}.  Then $\nu_{n,\sigma}$ is angle-monotone: for $\sigma>0$,
$\nu_{n,\sigma}=p\,\delta_1+(1-p)\,f\,m_n$ with $p=\Prob(1+A^2\not\succeq0)$ and
$f(t(\theta))=Z^{-1}F_{\sigma^2}(\theta)$, where $F_v$ is the function of
Lemma~\ref{lem:jacobian}(b) with $\epsilon=n-2m$ and $Z$ is a normalizing constant.  Consequently
the spectral gap of the walk on $\SO(n)$ with step distribution $\nu_{n,\sigma}$ is exactly
\[
  1-\lambda(\nu_{n,\sigma})=1-\frac1n\,\E[\tr R(A)],
\]
and $\lambda(\nu_{n,\sigma})$ is its least nonnegative $L^2$ contraction bound.
\end{corollary}

\begin{proof}
For $\sigma=0$, $\nu_{n,0}=\delta_1$.  Let $\sigma>0$, $B=\{A\in\mathfrak{so}(n):1+A^2\succ0\}$ and
$H=\{R\in\SO(n):\sym R\succ0\}$.  By Theorem~\ref{thm:principal-exists}(a) and
Propositions~\ref{prop:structure}(b) and~\ref{prop:principal}(b), $A\mapsto R(A)$ maps $B$
bijectively onto $H$, with inverse $R\mapsto\skw R$, and it is smooth on $B$, where
$R(A)=\sqrt{1+A^2}+A$ with $1+A^2\succ0$.  The boundary of $B$ is a null set, and $R(A)=1$ when
$1+A^2\not\succeq0$.  Give $\mathfrak{so}(n)$ the inner product $\ip XY=-\frac12\tr(XY)$, in which
the coordinates $A_{ij}$, $i<j$, are orthonormal and $A$ has the density
$\varphi(A)=(2\pi\sigma^2)^{-d/2}\exp(-\ip AA/2\sigma^2)$, and give $\SO(n)$ the left-invariant
metric equal to $\ip\cdot\cdot$ at $1$.  This metric is also right invariant, so its Riemannian
volume is a multiple $\kappa_nm_n$ of Haar measure.  In the orthonormal frames $(RJ_{ij})$ and
$(J_{ij})$, the differential of $R\mapsto\skw R$ at $R$ is $L_R(X)=\skw(RX)$.  By the change of
variables formula, for Borel $g\ge0$,
\[
  \E\bigl[g(R(A));A\in B\bigr]=\kappa_n\int_Hg(R)\,\varphi(\skw R)\,|\det L_R|\,dm_n(R).
\]
For $R=Qt(\theta)Q\transp$ we have $\skw R=Q(\skw t(\theta))Q\transp$ and
$L_R(X)=Q\,L_{t(\theta)}(Q\transp XQ)\,Q\transp$, so both factors depend only on $\theta$.  The map
$L_{t(\theta)}(X)=\frac12(t(\theta)X+Xt(\theta)\transp)$ preserves the decomposition of
$\mathfrak{so}(n)$ into the lines $\R E_k$, the four-dimensional spaces of matrices supported on
the blocks $(k,l)$ and $(l,k)$ for $k<l$, and, for odd $n$, the two-dimensional spaces supported
on the blocks $(k,n)$ and $(n,k)$; on them it acts by the three maps of
Lemma~\ref{lem:jacobian}(a).  Hence
\[
  \det L_{t(\theta)}=\prod_k\cos\theta_k\prod_{k<l}\Bigl(\frac{\cos\theta_k+\cos\theta_l}2\Bigr)^2
  \prod_k\Bigl(\frac{1+\cos\theta_k}2\Bigr)^{n-2m},
\]
and $\ip{\skw t(\theta)}{\skw t(\theta)}=\sum_k\sin^2\theta_k=\sum_k(1-\cos^2\theta_k)$.  On $H$
all $\cos\theta_k>0$.  So $\nu_{n,\sigma}$ has the stated form, with $f=0$ off $H$.  The density
$f$ is the product of the indicator of the closure of $H$ and a smooth function that vanishes on
the boundary of $H$, where some $\cos\theta_k=0$; such a product is Lipschitz.  It is invariant
under conjugation by $\Or(n)$, and $F_{\sigma^2}$ is nonincreasing on $[0,\pi]^m$ by
Lemma~\ref{lem:jacobian}(b).  Apply Theorem~\ref{thm:monotone}.
\end{proof}

With the upper bounds of Section~\ref{sec:son}, the spectral gap of the sampler with Gaussian
proposals is therefore known exactly for every $n$ and $\sigma$:
\[
  \gap=1-\frac{\E[\tr R(A)]}n=\frac{\E\fro{1-R(A)}^2}{2n},\qquad
  \frac{\E\bigl[\fro{A}^2;1+A^2\succeq0\bigr]}{2n}\le\gap\le(n-1)\sigma^2 .
\]

\begin{remark}[On Theorem~\ref{thm:monotone}]\label{rem:monotone}
\notformal\ (1) The proofs of Theorem~\ref{thm:monotone} and Corollary~\ref{cor:son-gauss} use the
Peter--Weyl theorem and Schur's lemma for $\SO(n)$, Rademacher's theorem, the normal form
$R=Qt(\theta)Q\transp$ of a rotation, and the change of variables formula on $\SO(n)$ together
with the identification of the Riemannian volume of a bi-invariant metric with a multiple of
Haar measure.  These are standard but are not in Mathlib.  Of them we have formalized only the
normal form, as a first step towards a formal proof: every $R\in\SO(n)$ is $Pt(\theta)P\transp$
with $P\in\SO(n)$ (\lname{exists_torusMat_conj} in \path{NormalForm.lean}).  The
computations that carry the two proofs are Lemmas~\ref{lem:grad} and~\ref{lem:jacobian}, which
are formalized.
(2) Lemma~\ref{lem:grad}(a) is where the standard representation wins.  The weights $w_\mu$ of
Step~6 form a probability distribution on the weights of the representation with
$\sum_\mu w_\mu\mu_k^2=c/d$, so $|\partial_k\tilde H|/c\le|\sin\theta_k|/d$, with equality for
$\psi$; the inequality $|\sin(j\theta)|\le|j|\,|\sin\theta|$ is the one that makes the first
Fourier mode dominate in Section~\ref{sec:sotwo}.  The comparison is between gradients, not
values.  Values of characters are not dominated pointwise, as the orbit measure of
Theorem~\ref{thm:so4} shows, but their gradients are, and an angle-monotone density only sees
gradients.
(3) For $n=3$ the theorem covers laws that are not supported where $\sym R\succeq0$, but not the
singular laws allowed in Theorem~\ref{thm:so3-gap}; for $n=2$ it contains the content of
Theorem~\ref{thm:unimodal}.
\end{remark}

\subsection{Numerical checks}

\begin{remark}[Numerical checks]\label{rem:numerics}
\notformal\ The script \path{numerics/check.py} accompanying this paper reports the following.
(1) For $n\in\{2,3,5,8\}$ the principal root is a rotation to machine precision, and the
fixed-point iteration $S\leftarrow S-\tfrac12\bigl((1+A+S)\transp(1+A+S)-1\bigr)$ started at
$S=0$ (Newton's method with the Jacobian frozen at the base point) converged to the principal
root in every trial in which it converged, as Theorem~\ref{thm:principal-exists}(c) says it
must.
(2) For $n=2$ and $\sigma\in\{0.05,0.2,0.5,1,3\}$, quadrature confirms
$|\hat\nu_\sigma(k)|\le\hat\nu_\sigma(1)$ for $2\le k\le200$ and the bound of
Theorem~\ref{thm:so2-size}.
(3) For $(n,\sigma)\in\{(3,0.2),(5,0.15),(8,0.1)\}$, a Monte Carlo estimate of
$1-\E[\tr R]/n$ agrees with one minus the lag-one autocorrelation of $X_{11}$ along a simulated
chain to within sampling error ($0.0423$ vs.\ $0.0427$; $0.0475$ vs.\ $0.0481$; $0.0364$
vs.\ $0.0360$), lies between the bounds of Proposition~\ref{prop:one-minus-lambda}, and is
close to $\tfrac12(n-1)\sigma^2$.
(4) For $n=3$ and $\sigma\in\{0.1,0.3,0.6,1.5\}$, $|\lambda_l|\le\lambda_1$ for $2\le l\le60$.
(5) For $n=4$ and $\sigma\in\{0.1,0.2,0.4,0.8\}$, among the irreducible representations with
highest weight parameters $j,k\le12$ the largest $|\lambda_\pi|$ is that of the standard
representation, as Corollary~\ref{cor:son-gauss} requires; the runners-up are the two
three-dimensional representations, on which the eigenvalue is $\lambda_2$.
(6) For the orbit measure $\nu_J$ of Theorem~\ref{thm:so4}, sampled with Haar-distributed
$Q\in\Or(4)$, the Monte Carlo estimates of $\E[\tr R]/4$ and of $\E[R_{11}R_{22}-R_{12}R_{21}]$
are $0.0000$ and $0.3331\pm0.0007$, against the exact values $0$ and $\tfrac13$.
(7) For $n\in\{4,5,6,7\}$ we tested Theorem~\ref{thm:monotone} on $96$ symmetric functions $F$ on
$[0,\pi]^m$ that are nonincreasing in each angle, other than the Gaussian density: indicators of
boxes, simplices and balls (which are not Lipschitz, but are limits of Lipschitz functions of the
same kind), $e^{-\kappa\sum_k\theta_k}$, $e^{-\kappa\max_k\theta_k}$,
$\prod_k(1+\cos\theta_k)^\kappa$, and random positive combinations of indicators of unions of
boxes.  The eigenvalue of the walk with density $F$ on each irreducible representation with
highest weight entries at most $7,6,4,4$ respectively, computed by quadrature from Weyl's
character and integration formulas, is at most $\lambda(\nu)$ in absolute value, and
$\lambda(\nu)\ge0$, up to rounding; for a density concentrated where every rotation angle equals
$2$, which is not monotone, $\lambda(\nu)<0$.
(8) For $(n,\sigma)\in\{(3,0.5),(4,0.35),(5,0.3),(6,0.25)\}$, the means of $\sum_k\cos\theta_k$
and of $\sum_k\cos2\theta_k$ over the accepted ones among $10^5$ simulated proposals agree with
their integrals against the density of Corollary~\ref{cor:son-gauss} to within $1.5$ standard
errors.
(9) On a grid of rotations all of whose angles lie in $(0,\pi/2)$, so that the symmetric part is
positive definite, the largest value of $|\chi_V|/\dim V-\tr R/n$ over the irreducible
representations $V$ other than the standard one with highest weight entries at most $w$ (for even
$n$, averaged over conjugation by $\Or(n)$) is positive for $n=4,6,8$ ($0.597$, $0.225$ and
$0.035$, with $w=8,5,3$) and is never positive for $n=3,5,7,9$ (with $w=60,18,8,3$, that is,
$59$, $188$, $163$ and $33$ representations).
\end{remark}

\section{Formalization}\label{sec:formal}

The Lean development is in the directory \path{lean/} next to this paper, as a Lake project
\path{SOnMCMC} depending on Mathlib.  It contains no \texttt{sorry}, no additional axioms and no
\texttt{native\_decide}; the command \texttt{\#print axioms} reports only \texttt{propext},
\texttt{Classical.choice} and \texttt{Quot.sound} for every declaration listed below.
Table~\ref{tab:files} lists the files and Table~\ref{tab:map} the correspondence between the
statements of the paper and the declarations, all in the namespace \texttt{SOnMCMC}.

\begin{table}[ht]
\caption{Files of the Lean development.}\label{tab:files}
\small
\begin{tabular}{@{}lp{0.63\textwidth}@{}}
\toprule
File & Contents\\
\midrule
\path{FiniteWalk.lean} & random walks on finite groups; Fourier transform; complete families\\
\path{Dihedral.lean} & representations, spectrum and spectral gap of the walk on $D_n$\\
\path{CompactWalk.lean} & Markov operator on $L^2$; stationarity, contraction, reversibility; gap\\
\path{Circle.lean} & Fourier multipliers and spectral gap on $\R/T\Z$\\
\path{CircleSampler.lean} & step angle; unimodal domination; the Gaussian sampler on the circle\\
\path{GapAsymptotics.lean} & Gaussian fourth moment; $|\gamma-\sigma^2/2|\le\frac32\sigma^4$; iterates\\
\path{Projection.lean} & algebra of the projection step; uniqueness of the principal root; step size\\
\path{PrincipalRoot.lean} & existence of the principal root; Newton iteration; the move of the sampler\\
\path{SO2.lean} & rotation matrices; the projection step in dimension two\\
\path{SOn.lean} & $\SO(n)$ as a compact group; eigenfunctions; gap bounds\\
\path{SO2Transport.lean} & transport along group isomorphisms; the sampler on the matrix group $\SO(2)$\\
\path{GaussianSampler.lean} & the Gaussian step distribution on $\SO(n)$; conjugation invariance; $n=2$\\
\path{SO3.lean} & character ratio inequality for $\SO(3)$\\
\path{SU2.lean} & $\mathrm{SU}(2)$ as a compact group; linear forms in the entries\\
\path{SU2Harmonic.lean} & polarization; the spaces $\mathcal H_k$; density\\
\path{ClassMeasure.lean} & class measures; extension by density; Pythagoras\\
\path{SU2Spectral.lean} & Schur's lemma for $\mathcal H_k$ by hand\\
\path{SU2Gap.lean} & eigenvalue recurrence; orthogonality; the $L^2$ bound on $\mathrm{SU}(2)$\\
\path{SO3Cover.lean} & the double cover $\mathrm{SU}(2)\to\SO(3)$\\
\path{SO3Gap.lean} & class decomposition; the spectral gap on $\SO(3)$\\
\path{SecondEigenvalue.lean} & two-by-two minors as eigenfunctions; the second upper bound\\
\path{SO4Counterexample.lean} & Haar measure on $\Or(n)$; orbit measures; the example on $\SO(4)$\\
\path{GradientBound.lean} & gradient bound on the torus; Jacobian of the projection; monotone angle density\\
\path{NormalForm.lean} & normal form $R=Pt(\theta)P\transp$ of a rotation (Remark~\ref{rem:monotone}; not cited by a numbered statement)\\
\bottomrule
\end{tabular}
\end{table}

Some points where the formal statements differ in form, but not in content, from the text:
\begin{itemize}
\item Functions in $L^p$ are functions $f$ with \texttt{MemLp f p m}, not equivalence classes;
$P_\nu f$ is defined pointwise by a Bochner integral, which is $0$ where the integrand is not
integrable (a null set of points $x$ when $f\in L^1$).
\item The spectral gap is expressed without choosing a value when it might not exist:
\texttt{IsSpectralGap m $\nu$ $\gamma$} says that $\gamma$ is the greatest Poincar\'e constant, and the
statements about contraction coefficients use \texttt{IsLeast}.
\item $\SO(n)$ is Mathlib's \texttt{Matrix.specialOrthogonalGroup (Fin n) $\mathbb{R}$}; $D_n$ is
Mathlib's \texttt{DihedralGroup n} with elements \texttt{r i}, \texttt{sr i}; the circle is
\texttt{AddCircle (2 * $\pi$)} and the Gaussian is \texttt{gaussianReal 0 v} with variance
$v=\sigma^2$ a nonnegative real.
\item The sign $\varepsilon=\pm1$ in Section~\ref{sec:dihedral} is encoded by a Boolean, and
$\cos(2\pi k/n)$ appears as the real part of $\omega^k$ (\path{psi_re} identifies the two).
\item In Theorem~\ref{thm:sampler} the random objects $A$ and $R$ are functions on an arbitrary
probability space, and Definition~\ref{def:sampler} is the hypothesis \texttt{hstep}.  In
Section~\ref{sec:construct} the probability space is $\R^{\{i<j\}}$ with the product of standard
Gaussian measures (\texttt{pairGauss}), the variance $v=\sigma^2$ is a nonnegative real, and
$R(A)$ is \texttt{principalStep A}, defined for every matrix $A$.
\item The move $M_g$ of Theorem~\ref{thm:zhg-move}(c) is \path{zhgMove}.  The projection it
uses is chosen by \path{Classical.choose}, which is harmless because the projection is unique.
\item $\mathrm{SU}(2)$ is Mathlib's \path{Matrix.specialUnitaryGroup} over $\C$ with index type
\texttt{Fin 2}.  Indices start at $0$, so $q_{11}$ is \texttt{q 0 0}.  The space $\mathcal H_k$
is \texttt{harm k}, a \texttt{Submodule} of the functions $\mathrm{SU}(2)\to\C$, and
$\lambda_k(\nu)$ is \texttt{su2Eig $\nu$ k}.  Invariance under inner automorphisms is
\path{IsConjInvariant}; invariance under $\Or(n)$, which is \path{ConjInvariant}, implies it on
$\SO(n)$.
\item Theorem~\ref{thm:so3-gap}(a) and Propositions~\ref{prop:su2-l2}(d) and~\ref{prop:classes}(b)
are stated for continuous functions; the passage to $L^2$ is made once, in
\path{so3_l2ContractionBound}.
\item $\lambda_2^{ij}(\nu)$ is \texttt{minorEig $\nu$ i j}, and the orbit measure $\nu_J$ is
\texttt{so4Step}.
\item In Lemma~\ref{lem:grad}, $H$ is \texttt{cosSum w $\mu$}, indexed by an arbitrary finite type,
and $\partial H/\partial\theta_k$ is \texttt{cosSumDeriv}, characterized by
\path{hasDerivAt_cosSum_update}.  In Lemma~\ref{lem:jacobian}, $R_a$ is \texttt{rotMat a},
$J$ is \texttt{skew2 1}, the third map is the $4\times4$ matrix \texttt{pairBlock a b} (identified
with the map by \path{pairBlock_mulVec}), $D_v$ and $F_v$ are \texttt{gaussCosDensity} and
\texttt{gaussAngleDensity} with a Boolean for $\epsilon$, and the statements hold for $v\ge0$.
\end{itemize}

\begin{longtable}{@{}p{0.27\textwidth}>{\raggedright\arraybackslash}p{0.69\textwidth}@{}}
\caption{Statements and their Lean declarations.}\label{tab:map}\\
\toprule
Statement & Lean declarations\\
\midrule
\endfirsthead
\toprule
Statement & Lean declarations\\
\midrule
\endhead
\bottomrule
\endfoot
Proposition \ref{prop:finite-basic} & \lname{walkOp_const}, \lname{sum_walkOp}, \lname{sum_walkOp_mul}, \lname{walkOp_walkOp}\\
Proposition \ref{prop:finite-fourier} & \lname{walkOp_matrixCoeff}, \lname{walkOp_eigenfunction}, \lname{fourierTransform_comm_of_class}, \lname{fourierTransform_conv}\\
Proposition \ref{prop:eigenbasis} & \lname{poincare_of_eigenbasis}, \lname{eigenvalue_mem_of_eigenbasis}\\
Proposition \ref{prop:compact-basic} & \lname{integrable_markovOp}, \lname{integral_markovOp}, \lname{memLp_markovOp_and_le}, \lname{integral_markovOp_mul}\\
Lemma \ref{lem:haar} & \lname{isMulRightInvariant_of_compact}\\
Proposition \ref{prop:compact-fourier} & \lname{markovOp_matrixCoeff}, \lname{markovOp_eigenfunction}, \lname{fourierMat_eq_conj}\\
Definition \ref{def:gap} & \lname{PoincareIneq}, \lname{IsSpectralGap}, \lname{L2ContractionBound}\\
Proposition \ref{prop:gap-tools} & \lname{PoincareIneq.le_of_eigenfunction}, \lname{L2ContractionBound.poincareIneq}, \lname{AddL2ContractionBound.iterate}\\
Lemma \ref{lem:dihedral-rep} & \lname{dihedralRep}, \lname{dihedralRep_r}, \lname{dihedralRep_sr}\\
Proposition \ref{prop:dihedral-fourier} & \lname{fourierTransform_dihedralRep}, \lname{psi_re}\\
Theorem \ref{thm:dihedral-spectrum} & \lname{dihedralEigFun_r}, \lname{dihedralEigFun_sr}, \lname{walkOp_dihedralEigFun}, \lname{dihedral_complete}, \lname{dihedral_spectrum}\\
Theorem \ref{thm:dihedral-gap} & \lname{dihedral_poincare}, \lname{dihedral_poincare_sharp}\\
Proposition \ref{prop:so2-model} & \lname{rotMat_mem}, \lname{rotMat_add}, \lname{exists_rotMat_eq}, \lname{rotMat_eq_one_iff}, \lname{abs_le_one_of_skewPart_eq}, \lname{principal_step_two}\\
Definition \ref{def:step} & \lname{stepAngle}, \lname{circleStep}, \lname{gaussStep}\\
Theorem \ref{thm:circle} & \lname{addMarkovOp_fourier}, \lname{fourierCoeff_addMarkovOp}, \lname{circle_l2ContractionBound}, \lname{circleChar_le_of_bound}, \lname{circle_isLeast_contraction}, \lname{circle_isSpectralGap}\\
Lemma \ref{lem:sin} & \lname{abs_sin_nat_mul_le}\\
Theorem \ref{thm:unimodal} & \lname{unimodal_integral_nonneg}\\
Theorem \ref{thm:so2-gap} & \lname{gaussLam}, \lname{cos_stepAngle}, \lname{circleChar_gaussStep}, \lname{circleChar_gaussStep_one}, \lname{gauss_char_le}, \lname{gauss_isSpectralGap}, \lname{gauss_isLeast_contraction}, \lname{gauss_iterate_bound}\\
Theorem \ref{thm:so2-size} & \lname{one_sub_gaussLam}, \lname{one_sub_gaussLam_ge}, \lname{one_sub_gaussLam_le}, \lname{integral_pow_four_gaussianReal}, \lname{abs_one_sub_gaussLam_sub_le}\\
Theorem \ref{thm:so2-matrix} & \lname{circleToSO_add}, \lname{circleToSO_bijective}, \lname{circleHomeo}, \lname{map_circleToSO_haar}, \lname{poincareIneq_map_iff}, \lname{l2ContractionBound_map_iff}, \lname{isSpectralGap_map_iff}, \lname{so2Step}, \lname{so2Step_eq_map}, \lname{so2_isSpectralGap}, \lname{so2_isLeast_contraction}\\
Proposition \ref{prop:structure} & \lname{exists_normal_iff}, \lname{symPart_mul_skewPart_comm}, \lname{symPart_sq}, \lname{add_mem_orthogonalGroup}, \lname{symPart_add_of_symm_of_skew}, \lname{skewPart_add_of_symm_of_skew}, \lname{skewPart_reverse}, \lname{symPart_reverse}, \lname{reverse_proposal_frobSq}\\
Proposition \ref{prop:principal} & \lname{det_eq_one_of_posSemidef}, \lname{principal_root_unique}, \lname{skewPart_conj}, \lname{symPart_conj}, \lname{posSemidef_symPart_conj}, \lname{frobSq_skewPart_le_card}\\
Proposition \ref{prop:step-size} & \lname{frobSq_one_sub}, \lname{frobSq_one_sub_eq}, \lname{frobSq_skewPart_le}, \lname{frobSq_one_sub_le}, \lname{card_sub_trace_bounds}\\
Proposition \ref{prop:son-haar} & \lname{SO.instIsTopologicalGroup}, \lname{SO.instCompactSpace}, \lname{haarSO}, \lname{haarSO.instIsHaarMeasure}, \lname{haarSO.instIsProbabilityMeasure}, \lname{haarSO.instIsMulRightInvariant}\\
Definition \ref{def:conj} & \lname{conjO}, \lname{ConjInvariant}, \lname{stdEig}\\
Theorem \ref{thm:entries} & \lname{fourierMat_std_eq}, \lname{markovOp_entry}, \lname{iterate_markovOp_entry}, \lname{integral_entry_haar}\\
Theorem \ref{thm:gap-upper} & \lname{gap_le_of_conjInvariant}\\
Proposition \ref{prop:one-minus-lambda} & \lname{one_sub_stdEig_eq_frob}, \lname{one_sub_stdEig_ge}, \lname{one_sub_stdEig_le_of_step}\\
Theorem \ref{thm:sampler} & \lname{integral_frobSq_of_var}, \lname{one_sub_stdEig_le_of_step}, \lname{gap_le_of_conjInvariant}, \lname{sampler_gap_le}\\
Theorem \ref{thm:principal-exists} & \lname{exists_principal_iff}, \lname{principalStep_spec}, \lname{eq_principalStep}, \lname{principalStep_mem}, \lname{posSemidef_symPart_principalStep}, \lname{principalStep_neg}, \lname{principalStep_conj}, \lname{one_add_newtonFrozen}, \lname{tendsto_visserS}, \lname{tendsto_visser}, \lname{tendsto_newtonFrozen}, \lname{posSemidef_of_tendsto_newtonFrozen}, \lname{isClosed_principal_domain}, \lname{measurable_principalStep}\\
Theorem \ref{thm:zhg-move} & \lname{IsPrincipalProjection}, \lname{isPrincipalProjection_iff}, \lname{isPrincipalProjection_reverse}, \lname{zhgMove}, \lname{zhgMove_eq}\\
Definition \ref{def:gauss-step} & \lname{Pairs}, \lname{skewOf}, \lname{pairGauss}, \lname{gaussTangent}, \lname{gaussStepFun}, \lname{gaussSOStep}\\
Theorem \ref{thm:gauss-sampler} & \lname{measurable_gaussStepFun}, \lname{gaussStepFun_spec}, \lname{integral_gaussTangent_sq}, \lname{pi_gaussian_map_of_isometry}, \lname{gaussSOStep_conjInvariant}, \lname{gaussSO_markovOp_entry}, \lname{gaussSO_gap_le_one_sub_stdEig}, \lname{gaussSO_gap_le}, \lname{principalStep_skew2}, \lname{gaussSOStep_two}, \lname{gaussSO_two_isSpectralGap}, \lname{gaussSO_two_isLeast_contraction}\\
Lemma \ref{lem:sin3} & \lname{abs_sin_mul_div_le}\\
Theorem \ref{thm:so3-char} & \lname{abs_so3CharRatio_le}, \lname{so3CharRatio_one_eq}, \lname{abs_so3CharRatio_le_one}\\
Corollary \ref{cor:so3} & \lname{so3_eigenvalue_dominated}\\
Proposition \ref{prop:su2-complete} & \lname{coeff_mem_of_forall_eval_mem}, \lname{harm}, \lname{lin_mul_lin_pow_mem}, \lname{polar_two}, \lname{lin_mul_lin_sub}, \lname{lin_mul_lin_pow_succ_mem}, \lname{harmAll}, \lname{mul_mem_harmAll}, \lname{conj_mem_harmAll}, \lname{mem_harmAll_of_mem_polyAlg}, \lname{polyAlg_topologicalClosure}, \lname{exists_harmAll_near}, \lname{continuous_of_mem_harm}, \lname{harm_parity}\\
Theorem \ref{thm:su2-schur} & \lname{su2Eig}, \lname{integral_mixed_eq_zero}, \lname{integral_first_column_pow}, \lname{exists_su2_first_column}, \lname{integral_lin_pow}, \lname{markovOp_lin_pow}, \lname{markovOp_of_mem_harm}\\
Theorem \ref{thm:su2-weyl} & \lname{integral_of_mem_harm}, \lname{su2Eig_trace_identity}, \lname{su2Eig_zero}, \lname{two_mul_su2Eig_one}, \lname{su2Eig_recurrence}, \lname{chebS}, \lname{su2Eig_mul_sin}, \lname{classMeasure}, \lname{classMeasure_isConjInvariant}, \lname{ae_classMeasure_trace}, \lname{norm_su2Eig_le_one}, \lname{norm_su2Eig_class_le}\\
Proposition \ref{prop:su2-l2} & \lname{classMeasure_map_inv}, \lname{integral_mul_conj_eq_zero_of_mem_harm}, \lname{integral_eq_zero_of_mem_harm}, \lname{integral_norm_sq_sum_of_orthogonal}, \lname{markovOp_sum_harm}, \lname{integral_norm_sq_markovOp_sum_harm}, \lname{exists_const_of_mem_harm_zero}, \lname{exists_even_harm_near}, \lname{integral_norm_sq_markovOp_le_of_approx}, \lname{su2_l2Bound_of_even}, \lname{su2_class_l2Bound}\\
Proposition \ref{prop:cover} & \lname{quatRot}, \lname{quatRot_mul}, \lname{quatRot_mul_transpose}, \lname{det_quatRot}, \lname{covMat_mul}, \lname{covMat_mem}, \lname{cover}, \lname{continuous_cover}, \lname{cover_negOne_mul}, \lname{covMat_torus}, \lname{covMat_yEl}, \lname{eq_rotX_of_entry}, \lname{cover_surjective}, \lname{measurePreserving_cover}, \lname{classMeasure_cover}, \lname{trace_covMat}, \lname{one_le_of_posSemidef_symPart}\\
Proposition \ref{prop:classes} & \lname{norm_integral_sq_le_weighted}, \lname{integral_classMeasure}, \lname{integral_norm_sq_markovOp_le_of_class}\\
Theorem \ref{thm:so3-gap} & \lname{one_le_trace_of_posSemidef}, \lname{markovOp_cover}, \lname{integral_comp_cover}, \lname{so3_class_l2Bound}, \lname{ConjInvariant.isConjInvariant}, \lname{so3_l2Bound_continuous}, \lname{one_third_le_stdEig}, \lname{so3_l2ContractionBound}, \lname{so3_isLeast_contraction}, \lname{so3_poincareIneq}, \lname{so3_isSpectralGap}\\
Corollary \ref{cor:so3-gauss} & \lname{ae_posSemidef_gaussSOStep}, \lname{gaussSO_three_isSpectralGap}, \lname{gaussSO_three_isLeast_contraction}\\
Theorem \ref{thm:second} & \lname{minor2}, \lname{minorEig}, \lname{integral_conjO}, \lname{integral_minor2}, \lname{markovOp_minor2}, \lname{minorEig_comm}, \lname{minorEig_swap}, \lname{minorEig_eq}, \lname{integral_minor2_haar}, \lname{integral_minor2_sq_pos}, \lname{gap_le_one_sub_minorEig}\\
Theorem \ref{thm:so4} & \lname{haarO}, \lname{orbitMeasure}, \lname{orbitMeasure_conjInvariant}, \lname{J4}, \lname{symPart_J4mat}, \lname{so4Step}, \lname{so4Step_conjInvariant}, \lname{ae_posSemidef_so4Step}, \lname{stdEig_so4Step}, \lname{e2_eq}, \lname{e2_conj_J4}, \lname{minorEig_so4Step}, \lname{so4_poincare_le}, \lname{so4_counterexample}\\
Lemma \ref{lem:grad} & \lname{abs_sin_int_mul_le}, \lname{cosProd}, \lname{cosSum}, \lname{cosSumDeriv}, \lname{cosProd_update}, \lname{hasDerivAt_cosSum_update}, \lname{abs_cosSumDeriv_le}, \lname{abs_sum_mul_le_of_sin}, \lname{gradient_domination}\\
Lemma \ref{lem:jacobian} & \lname{rotMat}, \lname{skew2}, \lname{skew_rotMat_mul_skew2}, \lname{det_half_rotMat_add_one}, \lname{pairBlock}, \lname{pairBlock_mulVec}, \lname{det_pairBlock}, \lname{gaussCosDensity}, \lname{gaussAngleDensity}, \lname{gaussCosDensity_nonneg}, \lname{gaussCosDensity_mono}, \lname{gaussAngleDensity_nonneg}, \lname{gaussAngleDensity_antitoneOn}\\
Theorem \ref{thm:monotone} & not formalized (Remark~\ref{rem:monotone})\\
Corollary \ref{cor:son-gauss} & not formalized (Remark~\ref{rem:monotone})\\
\end{longtable}

\subsection*{What is not formalized}
For convenience we collect the statements made in this paper that are \emph{not} machine
checked.
(1) Theorem~\ref{thm:monotone} and Corollary~\ref{cor:son-gauss}: the spectral gap of an
angle-monotone walk on $\SO(n)$, and in particular of the sampler with Gaussian proposals, is
exactly $1-\E[\tr R]/n$.  Their proofs are complete on paper, but they use the Peter--Weyl theorem,
Schur's lemma, Rademacher's theorem and the change of variables formula on $\SO(n)$, none of
which we have formalized, and the normal form of rotations, which we have
(Remark~\ref{rem:monotone}); the
computations they rest on, Lemmas~\ref{lem:grad} and~\ref{lem:jacobian}, are formalized.
(2) Remark~\ref{rem:model}: the relation between the chain $X_{t+1}=X_tR(A_{t+1})$ and a
particular implementation of the sampler of \cite{ZHG}.  The chain itself, the projection, the
reverse check, the Metropolis--Hastings test and the Newton iteration with frozen Jacobian are
formalized (Section~\ref{sec:construct}); the full Newton iteration with an updated Jacobian,
finite stopping rules and floating-point arithmetic are not.
(3) Remarks~\ref{rem:selfadjoint} and~\ref{rem:so3-classical}: the reformulation of
Definition~\ref{def:gap} through the spectral theorem, and the classical description of
$L^2(\SO(3))$ by the Peter--Weyl theorem.  Neither is used in a proof.
(4) The normalization in Remark~\ref{rem:so2-model}, the two heuristic remarks about
asymptotics (the leading-order size of the gap and the scaling of $\sigma$ with $n$), and all
numerical statements (Remarks~\ref{rem:so2-numerics} and~\ref{rem:numerics}).
(5) The prose between numbered statements, in particular the case $n=1$ after
Theorem~\ref{thm:dihedral-gap}, the interpretation of $\varrho$ by quaternions in
Section~\ref{sec:cover}, and the eigenvalue computation for rotations in two planes and the
discussion of even and odd $n$ after Theorem~\ref{thm:so4}.

\begin{thebibliography}{9}

\bibitem{Diaconis}
P.~Diaconis, \emph{Group Representations in Probability and Statistics}, Institute of
Mathematical Statistics Lecture Notes--Monograph Series, vol.~11, 1988.

\bibitem{Lean4}
L.~de~Moura and S.~Ullrich, \emph{The Lean~4 theorem prover and programming language}, in
Automated Deduction -- CADE~28, Lecture Notes in Computer Science, vol.~12699, Springer, 2021,
pp.~625--635.

\bibitem{Mathlib}
The mathlib Community, \emph{The Lean mathematical library}, in Proceedings of the 9th ACM
SIGPLAN International Conference on Certified Programs and Proofs (CPP 2020), 2020,
pp.~367--381.

\bibitem{Porod}
U.~Porod, \emph{The cut-off phenomenon for random reflections}, Ann. Probab. \textbf{24} (1996),
no.~1, 74--96.

\bibitem{Rosenthal}
J.~S.~Rosenthal, \emph{Random rotations: characters and random walks on $\mathrm{SO}(N)$}, Ann.
Probab. \textbf{22} (1994), no.~1, 398--423.

\bibitem{ZHG}
E.~Zappa, M.~Holmes-Cerfon, and J.~Goodman, \emph{Monte Carlo on manifolds: sampling densities
and integrating functions}, Comm. Pure Appl. Math. \textbf{71} (2018), no.~12, 2609--2647.

\end{thebibliography}

\end{document}
